% ==============================================================================
% Masterarbeit – Kapitel: Experimente und Ergebnisse (deutsche Fassung)
% Thema:  Dateneffiziente Personalisierung für webcam-basiertes Eye-Tracking
%         mithilfe vortrainierter Bildverarbeitungsmodelle
% Autor:  Khaled Rafei (MatNr. 930677)
% Datei:  experiments_results.tex (Ordner deutsch)
% Build:  python build_chapter.py experiments_results.tex
%         (benötigt references.bib und den Ordner figures/ neben dieser Datei)
% Daten:  Abschnitt 4.2 und der globale Verlustvergleich: exp_results_all.json
%         (exp.py --experiment all --sample-seed 0 --val-size 200, korrigierte
%         Metrik, 5 Wiederholungen, 98-%-KI, Lauf vom 22.10.2026).
%         Personalisierung (Abschnitte 4.4-4.6): personal_runs.jsonl
%         (personal_exp.py / personal_compare.ipynb, 24.10.2026).
% Hinweis: Inhalt, Labels und Zitierschlüssel entsprechen der englischen
%         Fassung F:\claude\@output\experiments_results.tex.
% ==============================================================================

\documentclass[
  11pt,
  a4paper,
  oneside,
  DIV=12,
  BCOR=8mm,
  parskip=half,
  headings=normal,
  bibliography=totoc,
  listof=totoc
]{scrreprt}

% ------------------------------------------------------------------------------
% Kodierung & Sprache
% ------------------------------------------------------------------------------
\usepackage[T1]{fontenc}
\usepackage[utf8]{inputenc}
\usepackage[english,ngerman]{babel}
\usepackage[autostyle=true,german=quotes]{csquotes}

% ------------------------------------------------------------------------------
% Schrift & Typografie
% ------------------------------------------------------------------------------
\usepackage{lmodern}
\usepackage{microtype}
\usepackage{setspace}
\setstretch{1.15}

% ------------------------------------------------------------------------------
% Mathematik & Einheiten
% ------------------------------------------------------------------------------
\usepackage{amsmath}
\usepackage{amssymb}
\usepackage{bm}
\usepackage{icomma}                  % Dezimalkomma im Mathemodus
\usepackage[locale=DE]{siunitx}

% ------------------------------------------------------------------------------
% Tabellen
% ------------------------------------------------------------------------------
\usepackage{booktabs}
\usepackage{tabularx}
\usepackage{array}
\usepackage{multirow}
\renewcommand{\arraystretch}{1.2}
\newcolumntype{Y}{>{\raggedright\arraybackslash}X}

% ------------------------------------------------------------------------------
% Abbildungen & Beschriftungen
% ------------------------------------------------------------------------------
\usepackage{graphicx}
\graphicspath{{figures/}}
\usepackage{float}
\usepackage[font=small,labelfont=bf]{caption}
\usepackage{subcaption}

% ------------------------------------------------------------------------------
% Listen, Farben, Kästen
% ------------------------------------------------------------------------------
\usepackage{enumitem}
\usepackage[dvipsnames,table]{xcolor}

% ------------------------------------------------------------------------------
% Literatur (natbib + bibtex: läuft überall, kein biber nötig)
% ------------------------------------------------------------------------------
\usepackage[numbers,sort&compress]{natbib}

% ------------------------------------------------------------------------------
% Links & Querverweise (hyperref zuletzt, cleveref danach)
% ------------------------------------------------------------------------------
\usepackage[hidelinks]{hyperref}
\usepackage[ngerman,capitalise,noabbrev]{cleveref}

% ------------------------------------------------------------------------------
% Eigene Befehle
% ------------------------------------------------------------------------------
\newcommand{\code}[1]{\texttt{#1}}
\setlength{\emergencystretch}{3em}   % vermeidet zu breite Zeilen
\newcommand{\Ediag}{E_{\mathrm{diag}}}
% Markierung der gewählten Option in den Hyperparametertabellen
\newcommand{\sel}{\ensuremath{^{\dagger}}}
% hervorgehobener Hinweiskasten
\newcommand{\notebox}[1]{%
  \par\noindent\fcolorbox{gray!60}{gray!8}{\parbox{\dimexpr0.97\linewidth-2\fboxsep-2\fboxrule}{\small #1}}\par}

% ==============================================================================
\begin{document}

\setcounter{chapter}{3}   % diese Datei ist Kapitel 4 der Arbeit
\chapter{Experimente und Ergebnisse}
\label{ch:results}

Dieses Kapitel berichtet über die Experimente auf dem GazeCapture-Datensatz
\cite{krafka2016itracker} und auf den eigenen Aufnahmesitzungen. Zuerst werden
das Modell, die Daten und das Evaluationsprotokoll beschrieben
(\cref{sec:setup}). Danach werden die Trainings-Hyperparameter
(\cref{sec:training}), die räumliche Verteilung des Vorhersagefehlers und die
Verlustfunktionen (\cref{sec:spatial}) sowie die Personalisierungssitzungen
(\cref{sec:sessions,sec:loading,sec:session_design}) analysiert.
\Cref{sec:discussion} diskutiert, was diese Ergebnisse für die weitere Arbeit
bedeuten.

Alle Werte wurden mit den Experimentprotokollen abgeglichen. Bei dieser
Prüfung wurde ein Fehler im Auswertungscode gefunden und korrigiert
(\cref{sec:metric}). Alle Modellfehler in diesem Kapitel werden mit der
korrigierten Metrik angegeben.

% ==============================================================================
\section{Versuchsaufbau}
\label{sec:setup}
% ==============================================================================

\subsection{Modell}
\label{sec:model}

\begin{itemize}[nosep]
  \item \textbf{Backbone:} ResNet-18 \cite{he2016resnet}, vortrainiert auf
        ImageNet (\code{ResNet18\_Weights.DEFAULT}).
  \item \textbf{Regressionskopf:} vier vollständig verbundene Schichten
        $512 \to 512 \to 256 \to 128 \to 2$ mit ReLU, Dropout $p=0,2$
        nach der zweiten Schicht und einer Ausgabe-Aktivierung (Standard:
        Sigmoid).
  \item \textbf{Optimierer:} AdamW, Weight Decay $10^{-5}$.
  \item \textbf{Standard-Verlust:} L1-Verlust (\code{nn.L1Loss}), siehe
        \cref{sec:losses}.
  \item \textbf{Early Stopping:} Geduld (Patience) von drei Epochen, höchstens
        500 Epochen.
\end{itemize}

\subsection{Daten und Splits}
\label{sec:data}

GazeCapture enthält in der hier verwendeten Version \num{347085} gelabelte
Frames von 262 Personen. Es werden zwei Splits verwendet:
\begin{itemize}[nosep]
  \item \code{norm\_subject}: personendisjunkter Split (keine Person kommt
        sowohl in den Trainings- als auch in den Testdaten vor);
  \item \code{norm\_random}: Split auf Frame-Ebene, dieselbe Person kann in
        Trainings- und Testdaten vorkommen.
\end{itemize}
Für die Hyperparameterstudie werden \num{1000} Trainings- und \num{200}
Testframes \emph{zufällig} aus den Split-Dateien gezogen (fester Seed 0,
\code{--sample-seed 0}). Für \code{norm\_subject} enthält die
Trainingsteilmenge 179 Personen und die Testteilmenge 45 Personen ohne
Überschneidung. (Frühere Läufe verwendeten die \emph{ersten} Zeilen der
sortierten Split-Dateien, d.\,h.\ nur fünf Trainingspersonen und eine
Testperson; diese Läufe werden für die Tabellen nicht mehr verwendet.)

\subsection{Evaluationsmetrik}
\label{sec:metric}

Alle Blickziele sind auf das Einheitsquadrat $[0,1]^2$ normiert
(\cref{sec:normalization}). Der Fehler ist die Wurzel aus dem mittleren
quadrierten \emph{euklidischen} Abstand zwischen Vorhersage $\hat{\bm p}_i$
und Ziel $\bm p_i$, angegeben in Prozent der Bildschirmdiagonale $\sqrt{2}$:
\begin{equation}
  \Ediag \;=\; 100\cdot
  \frac{\sqrt{\frac{1}{N}\sum_{i=1}^{N}\lVert\hat{\bm p}_i-\bm p_i\rVert_2^{2}}}{\sqrt{2}}
  \quad[\%].
  \label{eq:ediag}
\end{equation}
Ein Wert von \SI{0}{\percent} ist eine perfekte Vorhersage,
\SI{100}{\percent} ein Fehler so groß wie die gesamte Bildschirmdiagonale.


\subsection{Protokoll}
\label{sec:protocol}

Die Hyperparameterstudie wird mit dem Befehl
\begin{center}
  \code{python exp.py --experiment all --sample-seed 0 --val-size 200}
\end{center}
ausgeführt. Jeweils eine Achse (zum Beispiel die Batchgröße) wird verändert,
während alle anderen Werte auf der Standardkonfiguration bleiben
(\cref{tab:default_config}). Jede Konfiguration wird fünfmal trainiert. Das
Early Stopping verwendet eine personendisjunkte Validierungsmenge von 200
Frames, und angegeben wird der \emph{Test}fehler in der Epoche mit dem
niedrigsten Validierungsfehler; die Testmenge wird also nicht für die
Modellauswahl verwendet. Die Tabellen zeigen den mittleren Testfehler, das
zweiseitige \SI{98}{\percent}-Konfidenzintervall (KI) des Mittelwerts
(Student-$t$, vier Freiheitsgrade) und die mittlere Laufzeit (siehe unten).

\paragraph{Was die \enquote{Laufzeit} misst.} Die Spalte \emph{Laufzeit}
ist die tatsächlich verstrichene Zeit (Wall-Clock) \emph{eines vollständigen
Trainingslaufs}, gemittelt über die fünf Läufe. Ein Lauf beginnt mit der
Auswertung des untrainierten Modells (Epoche~0) und endet, wenn das Early
Stopping das Training beendet; er umfasst alle Trainingsepochen und die
Auswertung auf den Trainings-, Validierungs- und Testdaten nach jeder Epoche.
Sie ist \emph{nicht} die Zeit pro Epoche und enthält \emph{nicht} das Laden
der Bilder oder die Berechnung der zwischengespeicherten Backbone-Merkmale,
die einmal vor dem Lauf erfolgt. Da das Early Stopping die Läufe nach
unterschiedlich vielen Epochen beendet, spiegelt die Zeit auch wider, wie
lange sich eine Konfiguration noch verbessert.

\begin{table}[htbp]
  \centering
  \caption{Standardkonfiguration von \code{exp.py}. Jede Variation ändert
           genau einen dieser Werte.}
  \label{tab:default_config}
  \small
  \begin{tabular}{ll@{\hspace{2.5em}}ll}
    \toprule
    Parameter & Wert & Parameter & Wert \\
    \midrule
    Split               & \code{norm\_subject} & Trainierbare Schichten & nur FC \\
    Trainings-/Testgröße & \num{1000} / \num{200}\textsuperscript{a} & Ausgabe-Aktivierung & Sigmoid \\
    Batchgröße          & 32                   & Verlust          & L1 (\code{standard}) \\
    Lernrate            & $10^{-4}$            & Max.\ Epochen / Patience & 500 / 3 \\
    Weight Decay        & $10^{-5}$            & Validierungsmenge & 200 Frames\textsuperscript{b} \\
    Wiederholungen      & 5                    & & \\
    \bottomrule
  \end{tabular}\\[1mm]
  \parbox{0.9\linewidth}{\footnotesize \textsuperscript{a}\,zufällig gezogen,
  Seed 0. \textsuperscript{b}\,von einer Person, die weder in der Trainings-
  noch in der Testteilmenge vorkommt.}
\end{table}

\paragraph{Lesehilfe für die Tabellen.} In jeder Hyperparametertabelle ist
die gewählte Option mit $\dagger$ markiert. Die Spalte \enquote{$p$ vs.\
$\dagger$} gibt den $p$-Wert eines zweiseitigen Mann-Whitney-$U$-Tests
zwischen den fünf Läufen dieser Option und den fünf Läufen der gewählten
Option an. Ein Rangtest wird verwendet, weil fünf Läufe zu wenige sind, um
die Normalverteilung zu prüfen, und weil mehrere Optionen einzelne
Ausreißerläufe enthalten (z.\,B.\ ein ReLU-Lauf mit \SI{46,98}{\percent}).
Bei $5+5$ Läufen ist der kleinstmögliche $p$-Wert $0,008$; er bedeutet, dass
\emph{jeder} Lauf der einen Option besser ist als \emph{jeder} Lauf der
anderen. Ein Mittelwert ist nur dann \textbf{fett} gedruckt, wenn die
gewählte Option \emph{signifikant besser} ist als jede andere Option
derselben Tabelle ($p<0,05$ und niedrigerer Mittelwert).

Drei Eigenschaften dieses Protokolls sind zu beachten:
\begin{enumerate}[nosep]
  \item \textbf{Gleiche Daten in allen Wiederholungen.} Der Seed legt die
        Trainings-, Test- und Validierungsteilmengen fest, sodass sich die
        fünf Wiederholungen nur in Gewichtsinitialisierung, Datenreihenfolge
        und Dropout unterscheiden. Das KI misst daher die Zufälligkeit des
        Trainings, nicht die Variation zwischen verschiedenen
        Datenteilmengen.
  \item \textbf{Kleine Validierungsmenge.} Die 200 Validierungsframes stammen
        von einer einzigen Person. Das Early Stopping ist daher verrauscht,
        die Testmenge wird aber nie für die Modellauswahl verwendet.
  \item \textbf{Trainingszeit.} Bei Läufen mit \enquote{nur FC} wird der Kopf
        auf zwischengespeicherten ResNet-Merkmalen trainiert (512 Zahlen pro
        Bild statt des Bildes). Ein Lauf dauert daher nur etwa eine Sekunde,
        und die Unterschiede zwischen diesen Zeiten sind nicht relevant. Die
        Zeiten von FC + Layer~4 und vollem ResNet enthalten den Vorwärts- und
        Rückwärtsdurchlauf durch den Backbone in jeder Epoche und sind daher
        viel länger; sie sind mit den Zeiten von \enquote{nur FC} nicht
        direkt vergleichbar.
\end{enumerate}

% ------------------------------------------------------------------------------
\subsection{Normierung der Zielwerte}
\label{sec:normalization}

Die Roh-Blickziele von GazeCapture sind in Bildschirmpunkten des
Aufnahmegeräts angegeben, gemessen von der oberen linken Bildschirmecke. Ihr
Wertebereich hängt daher von der Bildschirmgröße und der Ausrichtung des
Geräts ab: Die \num{347085} Frames der 262 Personen stammen von zehn
Bildschirmformaten, von $320\times480$ bis $1024\times768$ Punkten,
aufgenommen im Hoch- und im Querformat. Die Ziele werden daher mit der
Bildschirmbreite $W$ und -höhe $H$ des jeweiligen Frames normiert
(\code{main.py}; Werte außerhalb von $[0,1]$ werden abgeschnitten, was
\SI{0,006}{\percent} der Frames betrifft):
\begin{equation}
  x_{\mathrm{norm}} = \frac{x}{W},\qquad y_{\mathrm{norm}} = \frac{y}{H}.
\end{equation}
\Cref{fig:normalization} zeigt die Wirkung. Ohne Normierung konzentrieren
sich die Ziele der vielen Smartphone-Frames im oberen linken Bereich, und
dieselbe Koordinate bedeutet auf jedem Gerät eine andere Bildschirmposition
(\cref{fig:raw_targets}); die Fehlerwerte sind daher zwischen Geräten nicht
vergleichbar (linke Spalte). Nach der Normierung liegen alle Ziele in
$[0,1]^2$, und \cref{eq:ediag} ist wohldefiniert (rechte Spalte). Alle
weiteren Experimente verwenden normierte Ziele.

\begin{figure}[htbp]
  \centering
  \begin{subfigure}[t]{0.43\textwidth}
    \centering
    \includegraphics[width=\textwidth]{raw_targets.png}\\[2mm]
    \includegraphics[width=\textwidth]{raw_learning_curve.jpg}
    \caption{Roh-Ziele in Bildschirmpunkten (Dichte aller Frames,
             logarithmische Farbskala); gestrichelt: Umrisse der
             Bildschirmformate. Unten: Lernkurve eines frühen Laufs auf
             Rohkoordinaten; die Werte sind keine Prozente einer gemeinsamen
             Bildschirmdiagonale und mit den anderen Abbildungen nicht
             vergleichbar.}
    \label{fig:raw_targets}
  \end{subfigure}
  \hfill
  \begin{subfigure}[t]{0.52\textwidth}
    \centering
    \includegraphics[width=0.82\textwidth]{normalized_targets.png}\\[2mm]
    \includegraphics[width=\textwidth]{norm_test_train_const_rand.png}
    \caption{Normierte Ziele (gleiche Darstellung; gestrichelt:
             Einheitsquadrat). Die hellen Zellen sind Kalibrierpunkte, die
             auf allen Geräten an denselben relativen Positionen gezeigt
             werden. Unten: Lernkurve auf \code{norm\_random}
             (\num{1000}/\num{200}) mit der Metrik aus \cref{eq:ediag}.
             Epoche~0 ist das untrainierte Modell. Die waagerechten Linien
             sind Baselines des \emph{gesamten} Datensatzes; siehe Text.}
    \label{fig:normalized_targets}
  \end{subfigure}
  \caption{Roh-Blickziele im Vergleich zu normierten Blickzielen.}
  \label{fig:normalization}
\end{figure}

\paragraph{Warum beginnt die Kurve bei der Konstant- und nicht bei der Zufalls-Baseline?}
In \cref{fig:normalized_targets} beginnen die Kurven in Epoche~0, also
\emph{vor} dem ersten Trainingsschritt, bei \SI{26,77}{\percent}. Man könnte
erwarten, dass sich das untrainierte Modell wie die Zufalls-Baseline (etwa
\SI{40}{\percent}) verhält. Diese Erwartung trifft aus folgendem Grund nicht
zu.

Die Zufalls-Baseline rät für jedes Bild einen \emph{anderen}, gleichverteilten
Punkt über den ganzen Bildschirm. Ein untrainiertes Netz tut das nicht.
Seine Gewichte sind mit kleinen Zufallswerten initialisiert, daher erzeugt
die letzte lineare Schicht für jedes Bild Werte nahe $0$, und die Sigmoid
bildet sie auf $\sigma(0)=0,5$ ab. Das untrainierte Modell sagt daher für
alle Bilder fast \emph{denselben} Punkt vorher, die Bildschirmmitte
$(0{,}5;\,0{,}5)$. Genau das tut die Konstant-Baseline. Die richtige
Erwartung lautet daher:
\begin{itemize}[nosep]
  \item Epoche~0 $\approx$ Konstant-Baseline \emph{derselben} Testdaten
        (\enquote{immer in die Mitte schauen});
  \item das Training muss die Kurve \emph{unter} diesen Wert bringen; erst
        dann hat das Modell etwas aus den Bildern gelernt;
  \item die Zufalls-Baseline ist eine obere Referenz für \enquote{reines
        Raten}; ein Modell sollte nie darüber liegen.
\end{itemize}
Die gemessenen Werte bestätigen das: Auf den 200 Testbildern dieses Laufs
beträgt die Konstant-Baseline \SI{26,76}{\percent} (Trainingsmittelwert) bzw.\
\SI{26,71}{\percent} (Bildschirmmitte), was dem Startwert von
\SI{26,77}{\percent} entspricht. Die waagerechte Linie in der alten Abbildung
(\SI{27,25}{\percent}) wurde auf dem \emph{gesamten} Datensatz berechnet und
liegt daher etwas höher; dadurch schien die Kurve \enquote{unterhalb der
Baseline} zu beginnen. Der Plot-Code wurde korrigiert: Die Baselines werden
jetzt auf derselben Trainings-/Testteilmenge wie die Kurven berechnet
(\code{metrics.subset\_baselines}; \code{Compare\_notebook.ipynb},
Abschnitt~4, der diese Abbildung auch neu erzeugt).

% ------------------------------------------------------------------------------
\subsection{Baselines}
\label{sec:baselines}

Es werden zwei Referenzschätzer verwendet. Die \emph{Konstant-Baseline} sagt
für jedes Testbild den mittleren Zielwert der Trainingsdaten vorher. Die
\emph{Zufalls-Baseline} zieht Vorhersagen gleichverteilt aus $[0,1]^2$. Ein
Modell ist nur dann nützlich, wenn es deutlich besser ist als die
Konstant-Baseline.

\begin{table}[htbp]
  \centering
  \caption{Baseline-Fehler $\Ediag$ in \si{\percent}. Oben: gesamter
           Datensatz (\code{exp.py}, Bootstrap mit 30 Stichproben zu je
           2000 Frames, Mittelwert und \SI{98}{\percent}-KI). Unten: dieselben
           zufälligen Trainings-/Testteilmengen (Seed 0), die für das
           Modelltraining verwendet werden (kein KI, eine feste Teilmenge).}
  \label{tab:baselines}
  \small
  \begin{tabular}{lcc}
    \toprule
    Daten & Konstant-Baseline & Zufalls-Baseline \\
    \midrule
    Gesamter Datensatz (\num{347085} Frames) & 27,28 \,[27,15; 27,40] & 39,82 \,[39,60; 40,04] \\
    \midrule
    \code{norm\_subject}, \num{1000}/\num{200}   & 28,46 & 40,36 \\
    \code{norm\_subject}, \num{10000}/\num{2000} & 27,62 & 39,96 \\
    \code{norm\_random}, \num{1000}/\num{200}    & 27,66 & 39,83 \\
    \bottomrule
  \end{tabular}
\end{table}

Die Baselines des früheren Entwurfs (\SI{30,99}{\percent} und
\SI{42,70}{\percent}) stammten aus einem älteren Notebook-Lauf, der nicht
Teil der Protokolle ist; sie sind durch die protokollierten Werte in
\cref{tab:baselines} ersetzt. Der untere Teil der Tabelle ist die richtige
Referenz für alle Modellergebnisse in \cref{sec:training}.

% ==============================================================================
\section{Trainingskonfiguration}
\label{sec:training}
% ==============================================================================

Alle Ergebnisse dieses Abschnitts stammen aus einem Lauf
(\code{exp\_results\_all.json}; \code{exp.py}, korrigierte Metrik, zufällige
Teilmengen mit Seed 0, validierungsbasiertes Early Stopping, fünf
Wiederholungen pro Konfiguration). Die Konstant-Baseline der
Standardkonfiguration auf derselben Teilmenge beträgt \SI{28,46}{\percent}.

% ------------------------------------------------------------------------------
\subsection{Batchgröße}
\label{sec:batch}

\begin{table}[htbp]
  \centering
  \caption{Batchgröße. Mittlerer Testfehler über fünf Läufe mit
           \SI{98}{\percent}-KI und mittlere Dauer eines vollständigen
           Trainingslaufs (Kopf auf zwischengespeicherten Merkmalen, siehe
           \cref{sec:protocol}).}
  \label{tab:batch_size}
  \small
  \begin{tabular}{ccccc}
    \toprule
    Batchgröße & Mittl.\ Testfehler (\%) & \SI{98}{\percent}-KI (\%) & $p$ vs.\ $\dagger$ & Laufzeit (min) \\
    \midrule
    32\sel{}     & 28,39 & [28,21; 28,58] & ---  & 0,019 \\
    64           & 28,35 & [28,29; 28,41] & 0,69 & 0,012 \\
    128          & 28,38 & [28,30; 28,47] & 0,84 & 0,009 \\
    \bottomrule
  \end{tabular}
\end{table}

\Cref{tab:batch_size} zeigt, dass die drei Batchgrößen denselben Testfehler
ergeben: Die Mittelwerte unterscheiden sich um höchstens \num{0,05}
Prozentpunkte, die Konfidenzintervalle überlappen, und der Rangtest findet
keinen Unterschied ($p\ge0,69$). Die Trainingszeiten liegen unter einer
Sekunde pro Lauf, weil nur der kleine Kopf auf zwischengespeicherten Merkmalen
trainiert wird; die Zeitunterschiede sind daher nicht relevant.

Eine Batchgröße von 32 wird als Standard beibehalten. Das Experiment zeigt
keinen Nachteil dieses Werts, und für eine kleine Batchgröße im
Personalisierungssetting sprechen zwei praktische Argumente: (i) Eine
persönliche Sitzung enthält nur einige hundert Trainingsbilder, sodass eine
kleine Batchgröße mehr Parameterupdates pro Epoche ergibt, und (ii) dem
Gradientenrauschen kleiner Batches wird üblicherweise eine leicht
regularisierende Wirkung zugeschrieben. Die Daten dieses Experiments
\emph{belegen} jedoch keinen Vorteil von 32 gegenüber 64 oder 128.

% ------------------------------------------------------------------------------
\subsection{Lernrate}
\label{sec:lr}

\begin{table}[htbp]
  \centering
  \caption{Lernrate. Mittlerer Testfehler über fünf Läufe mit
           \SI{98}{\percent}-KI und mittlere Dauer eines vollständigen
           Trainingslaufs.}
  \label{tab:learning_rate}
  \small
  \begin{tabular}{ccccc}
    \toprule
    Lernrate & Mittl.\ Testfehler (\%) & \SI{98}{\percent}-KI (\%) & $p$ vs.\ $\dagger$ & Laufzeit (min) \\
    \midrule
    $10^{-3}$          & 28,43 & [28,33; 28,54] & 0,15  & 0,018 \\
    $10^{-4}$\sel{}    & 28,37 & [28,32; 28,42] & ---   & 0,020 \\
    $10^{-5}$          & 28,33 & [28,28; 28,38] & 0,056 & 0,022 \\
    $10^{-6}$          & 28,37 & [28,21; 28,52] & 0,42  & 0,055 \\
    \bottomrule
  \end{tabular}
\end{table}

\begin{figure}[htbp]
  \centering
  \includegraphics[width=0.78\textwidth]{Test-Error_diff_learning-rate.png}
  \caption{Testfehlerkurven einzelner Trainingsläufe mit vier Lernraten
           (\code{norm\_random}, \num{1000}/\num{200}, 15 Epochen ohne Early
           Stopping, Metrik aus \cref{eq:ediag}). Der erste Punkt (Epoche~1
           auf der Achse) ist das untrainierte Modell.}
  \label{fig:lr_curves}
\end{figure}

\Cref{tab:learning_rate} zeigt in der Variation mit fünf Läufen keinen
signifikanten Unterschied zwischen den vier Lernraten: Alle Mittelwerte liegen
zwischen \SI{28,33}{\percent} und \SI{28,43}{\percent}, die
Konfidenzintervalle überlappen, und kein Rangtest ist signifikant
($p\ge0,056$). Der Endfehler allein kann die Wahl daher nicht begründen. Die
Wahl von $\eta = 10^{-4}$ beruht auf dem \emph{Konvergenzverhalten} in
\cref{fig:lr_curves}:
\begin{itemize}[nosep]
  \item $\bm{\eta = 10^{-3}}$: Die Schritte sind zu groß. Nach einem schnellen
        ersten Abfall schwankt der Testfehler stark und steigt wieder an
        (Spitzen von etwa \SI{27,1}{\percent} und \SI{28,4}{\percent}); der Lauf
        konvergiert nicht zu einem stabilen Minimum.
  \item $\bm{\eta = 10^{-4}}$: erreicht den niedrigsten Testfehler aller
        Kurven (etwa \SI{23,6}{\percent} um Epoche~10). Die Kurve schwankt
        noch, aber viel weniger als bei $10^{-3}$, und das Early Stopping
        kann den besten Zustand behalten.
  \item $\bm{\eta = 10^{-5}}$ \textbf{und} $\bm{10^{-6}}$: Die Kurven sind sehr
        glatt und konvergieren monoton, bleiben aber auf einem Plateau von
        etwa \SI{24,3}{\percent}. Die Schrittweite ist zu klein, um diesen
        Bereich innerhalb des Trainingsbudgets zu verlassen; das Ergebnis ist
        stabil, aber deutlich nicht das beste erreichbare.
\end{itemize}
Die Trainingszeit wird hier \emph{nicht} als Kriterium verwendet: Sie liegt
für alle Lernraten unter einer Sekunde (zwischengespeicherte Merkmale), und
die Entscheidung hängt nur vom Konvergenzverhalten und vom erreichbaren
Fehler ab. Zu beachten ist, dass \cref{fig:lr_curves} einzelne Läufe zeigt;
die Variation mit fünf Läufen in \cref{tab:learning_rate} bestätigt keinen
signifikanten Unterschied.

% ------------------------------------------------------------------------------
\subsection{Trainierbare Schichten}
\label{sec:layers}

Es werden drei Fine-Tuning-Tiefen verglichen (\cref{tab:params}): nur der
Regressionskopf (nur FC), der Kopf plus die letzte Residualstufe
(\code{layer4}) und das gesamte Netz.

\begin{table}[htbp]
  \centering
  \caption{Trainierbare Parameter der drei Fine-Tuning-Konfigurationen
           (ResNet-18 mit Regressionskopf, insgesamt \num{11603650}
           Parameter).}
  \label{tab:params}
  \small
  \begin{tabular}{lrrr}
    \toprule
    & Nur FC & FC + Layer 4 & Volles ResNet \\
    \midrule
    Trainierbare Parameter         & \num{427138} & \num{8820866} & \num{11603650} \\
    Anteil an allen Parametern     & \SI{3,68}{\percent} & \SI{76,02}{\percent} & \SI{100,00}{\percent} \\
    Backbone-Merkmale speicherbar  & ja & nein & nein \\
    \bottomrule
  \end{tabular}
\end{table}

\begin{table}[htbp]
  \centering
  \caption{Trainierbare Schichten. Mittlerer Testfehler über fünf Läufe mit
           \SI{98}{\percent}-KI und mittlere Dauer eines vollständigen
           Trainingslaufs. Die Zeit für \enquote{nur FC} enthält keinen
           Vorwärtsdurchlauf durch den Backbone (zwischengespeicherte
           Merkmale).}
  \label{tab:layers}
  \small
  \begin{tabular}{lcccc}
    \toprule
    Trainierbare Schichten & Mittl.\ Testfehler (\%) & \SI{98}{\percent}-KI (\%) & $p$ vs.\ $\dagger$ & Laufzeit (min) \\
    \midrule
    Nur FC\sel{}           & 28,36 & [28,23; 28,49] & ---  & 0,015 \\
    FC + Layer 4           & 29,07 & [26,58; 31,55] & 0,22 & 9,160 \\
    Volles ResNet          & 29,39 & [26,99; 31,78] & 0,31 & 15,366 \\
    \bottomrule
  \end{tabular}
\end{table}

\enquote{Nur FC} hat den niedrigsten mittleren Fehler und mit Abstand das
schmalste Intervall (\cref{tab:layers}). Das Trainieren weiterer Schichten
hilft nicht; es macht das Ergebnis \emph{weniger stabil}: Einer der fünf
Läufe mit FC + Layer~4 (\SI{31,72}{\percent}) und zwei der fünf Läufe mit
vollständigem Fine-Tuning (\SI{30,70}{\percent} und \SI{31,18}{\percent})
enden deutlich über der Konstant-Baseline (\SI{28,46}{\percent}), während alle
Läufe mit \enquote{nur FC} zwischen \SI{28,28}{\percent} und
\SI{28,47}{\percent} bleiben. Mit nur fünf Läufen sind die Unterschiede nicht
signifikant ($p=0,22$ und $0,31$), daher ist der Wert für \enquote{nur FC}
nicht fett gedruckt. Die Trainingszeit wächst stark mit der Zahl der
trainierbaren Parameter: \SI{9,2}{\minute} pro Lauf für FC + Layer~4 und
\SI{15,4}{\minute} für das gesamte Netz.

\enquote{Nur FC} wird aus drei Gründen für die Personalisierungsexperimente
gewählt:
\begin{enumerate}[nosep]
  \item \textbf{Kein Genauigkeitsverlust, mehr Stabilität:} In diesem
        Experiment senkt das Trainieren weiterer Schichten den Fehler nicht,
        sondern erhöht die Streuung zwischen den Läufen.
  \item \textbf{Robustheit bei wenigen Daten:} Eine persönliche Sitzung hat
        nur einige hundert Bilder. Millionen von Backbone-Parametern mit so
        wenigen Beispielen zu aktualisieren, birgt die Gefahr der
        Überanpassung und verzerrt die vortrainierten Merkmale;
        \citet{kumar2022finetuning} zeigen, dass vollständiges Fine-Tuning
        vortrainierte Repräsentationen beschädigen und schlechter
        generalisieren kann als das Training nur des Kopfes.
  \item \textbf{Aufwand:} Nur \SI{3,68}{\percent} der Parameter werden
        trainiert, und die Backbone-Merkmale können einmal berechnet und
        wiederverwendet werden. Das macht die Personalisierung schnell und
        speichersparend.
\end{enumerate}

% ------------------------------------------------------------------------------
\subsection{Ausgabe-Aktivierung}
\label{sec:activation}

Die Ausgabe-Aktivierung wird nach der letzten linearen Schicht angewendet,
die $(x,y)$ erzeugt. Sigmoid bildet die Ausgabe auf $(0,1)$ ab, was dem
normierten Koordinatenbereich entspricht. Die Gauß-Aktivierung
$g(z)=\exp(-z^2/2\sigma^2)$ mit $\sigma=1$ bildet ebenfalls auf $(0,1]$ ab, ist
aber nicht monoton. ReLU entfernt nur negative Werte, und \enquote{Keine}
lässt die lineare Ausgabe unverändert.

\begin{table}[htbp]
  \centering
  \caption{Ausgabe-Aktivierung. Mittlerer Testfehler über fünf Läufe mit
           \SI{98}{\percent}-KI und mittlere Dauer eines vollständigen
           Trainingslaufs.}
  \label{tab:activation}
  \small
  \begin{tabular}{lcccc}
    \toprule
    Aktivierung & Mittl.\ Testfehler (\%) & \SI{98}{\percent}-KI (\%) & $p$ vs.\ $\dagger$ & Laufzeit (min) \\
    \midrule
    Sigmoid\sel{}    & \textbf{28,36} & [28,29; 28,43] & ---   & 0,022 \\
    Gauß             & 31,29 & [28,44; 34,14] & 0,008 & 0,019 \\
    ReLU             & 32,90 & [19,69; 46,10] & 0,008 & 0,028 \\
    Keine            & 29,39 & [28,42; 30,36] & 0,008 & 0,032 \\
    \bottomrule
  \end{tabular}
\end{table}

Sigmoid ist der einzige klare Gewinner der Hyperparameterstudie
(\cref{tab:activation}): Sie hat den niedrigsten mittleren Fehler und das
schmalste Intervall, und \emph{jeder} Sigmoid-Lauf ist besser als \emph{jeder}
Lauf der drei anderen Aktivierungen ($p=0,008$ in allen drei Vergleichen).
Der Wert ist daher fett gedruckt. Die anderen Aktivierungen sind auch weniger
stabil: Ein ReLU-Lauf endet bei \SI{46,98}{\percent} und damit schlechter als
die Zufalls-Baseline. Ein plausibler Grund ist, dass ReLU-Ausgaben für alle
Eingaben null werden können (\enquote{tote} Ausgaben), wovon sich der Gradient
nicht erholt. Ohne begrenzende Aktivierung (Keine, ReLU) kann das Netz
außerdem Punkte außerhalb des Bildschirms vorhersagen. Sigmoid wird gewählt,
weil sie die Ausgabe auf den gültigen Bildschirmbereich begrenzt und den
niedrigsten und stabilsten Fehler ergab.

% ------------------------------------------------------------------------------
\subsection{Datensatzgröße}
\label{sec:size}

\begin{table}[htbp]
  \centering
  \caption{Datensatzgröße (Training/Test). Mittlerer Testfehler über fünf
           Läufe mit \SI{98}{\percent}-KI, mittlere Dauer eines vollständigen
           Trainingslaufs und Konstant-Baseline der jeweiligen Testteilmenge.}
  \label{tab:dataset_size}
  \small
  \setlength{\tabcolsep}{3.5pt}
  \begin{tabular}{lccccc}
    \toprule
    Training / Test & Mittl.\ Testfehler (\%) & \SI{98}{\percent}-KI (\%) & $p$ vs.\ $\dagger$ & Laufzeit (min) & Konst.-Baseline (\%) \\
    \midrule
    \num{10000} / \num{2000}      & 27,64 & [27,55; 27,73] & 0,008 & 0,167 & 27,62 \\
    \num{1000} / \num{200}\sel{}  & 28,39 & [28,33; 28,44] & ---   & 0,018 & 28,46 \\
    \bottomrule
  \end{tabular}
\end{table}

Die größere Konfiguration hat einen niedrigeren mittleren Fehler
(\cref{tab:dataset_size}). Die beiden Konfigurationen werden jedoch auf
\emph{unterschiedlichen} Testmengen ausgewertet, und jeder Mittelwert ist
praktisch gleich der Konstant-Baseline seiner eigenen Testmenge
(\SI{27,64}{\percent} vs.\ \SI{27,62}{\percent} und \SI{28,39}{\percent} vs.\
\SI{28,46}{\percent}). Der Unterschied von \num{0,75} Prozentpunkten wird
daher durch die unterschiedlichen Testdaten verursacht und nicht durch mehr
Trainingsdaten, auch wenn der Rangtest signifikant ist ($p=0,008$). Die kleine
Konfiguration (\num{1000}/\num{200}) wird für die Variationen beibehalten,
weil sie etwa zehnmal schneller ist und der Datenmenge einer
Personalisierungssitzung entspricht.

% ------------------------------------------------------------------------------
\subsection{Split-Typ}
\label{sec:split}

\begin{table}[htbp]
  \centering
  \caption{Split-Typ. Mittlerer Testfehler über fünf Läufe mit
           \SI{98}{\percent}-KI und Konstant-Baseline der jeweiligen
           Testteilmenge.}
  \label{tab:split}
  \small
  \begin{tabular}{lcccc}
    \toprule
    Split & Mittl.\ Testfehler (\%) & \SI{98}{\percent}-KI (\%) & $p$ vs.\ $\dagger$ & Konstant-Baseline (\%) \\
    \midrule
    \code{norm\_random}           & 27,72 & [27,64; 27,80] & 0,008 & 27,66 \\
    \code{norm\_subject}\sel{}    & 28,39 & [28,26; 28,53] & ---   & 28,46 \\
    \bottomrule
  \end{tabular}
\end{table}

Der Zufallssplit ergibt den niedrigeren Fehler (\cref{tab:split}). Auch hier
stammt der größte Teil des Unterschieds aus den Testteilmengen: Die
Konstant-Baseline ist auf der \code{norm\_random}-Teilmenge bereits um
\num{0,80} Prozentpunkte niedriger, und beide Modelle liegen innerhalb von
\num{0,07} Prozentpunkten ihrer eigenen Baseline. Zusätzlich erlaubt der
Zufallssplit, dass dieselben Personen in Trainings- und Testdaten vorkommen,
was die Aufgabe erleichtert und die Schätzung optimistisch macht. Der
personendisjunkte Split \code{norm\_subject} wird daher als Standard
verwendet: Er ist der einzige, der die Generalisierung auf unbekannte
Personen misst, also die Situation vor der Personalisierung.

% ==============================================================================
\section{Räumliche Fehleranalyse und Verlustfunktionen}
\label{sec:spatial}
% ==============================================================================

Der mittlere Fehler aus \cref{sec:training} verbirgt, \emph{wo} auf dem
Bildschirm die Fehler auftreten. Dieser Abschnitt analysiert zuerst die
räumliche Fehlerverteilung (\cref{sec:partition,sec:heatmaps}). Das
beobachtete Muster motivierte positionsgewichtete Verlustfunktionen
(\cref{sec:losses}), die in \cref{sec:loss_results} ausgewertet werden.

\subsection{Bildschirmeinteilung}
\label{sec:partition}

Der normierte Bildschirm $[0,1]^2$ wird durch ein $3\times3$-Referenzraster
mit drei Regionen beschrieben:
\begin{align*}
  \text{Mitte:}  &\quad (0,50;\,0,50),\\
  \text{Ränder:} &\quad (0,10;\,0,50),\ (0,50;\,0,10),\ (0,90;\,0,50),\ (0,50;\,0,90),\\
  \text{Ecken:}  &\quad (0,10;\,0,10),\ (0,10;\,0,90),\ (0,90;\,0,10),\ (0,90;\,0,90).
\end{align*}
Ein Testbeispiel mit Ziel $\bm p$ gehört zum Rasterpunkt $\bm g$, wenn
$\lVert\bm p-\bm g\rVert_2 \le r$ mit $r = 0,15$ gilt. Der Fehler eines
Rasterpunkts ist der mittlere euklidische Fehler $\lVert\hat{\bm p}-\bm
p\rVert_2$ seiner Beispiele (in normierten Einheiten, d.\,h.\ als Anteil der
Bildschirmseite; die Diagonale ist $\sqrt2 \approx 1,41$). Der Fehler einer
Region ist der Mittelwert über alle Beispiele ihrer Rasterpunkte. Diese
Definitionen werden in allen Abbildungen und Tabellen dieses Abschnitts
verwendet.

\subsection{Fehler-Heatmaps}
\label{sec:heatmaps}

\Cref{fig:heatmaps} verwendet die 200 Testvorhersagen des Standardmodells
(\code{norm\_subject}, L1-Verlust; Datei
\code{results/heatmap/base\_loss1.csv}). Für eine Cursorposition $\bm c$
(Mitte, Rand oder Ecke) wird der mittlere Fehler $\bar e(\bm c)$ aller
Beispiele innerhalb des Radius $r$ berechnet und als isotrope Gaußfunktion
dargestellt:
\begin{equation}
  H(x,y) = \bar e(\bm c)\,
  \exp\!\left(-\frac{\lVert (x,y)-\bm c\rVert_2^2}{2\sigma^2}\right),
  \qquad \sigma = 0,20 .
\end{equation}
Die Farbe kodiert daher die \emph{Größe} des lokalen mittleren Fehlers, und
die Lage des hellen Flecks zeigt, \emph{wo} er gemessen wurde. Alle drei
Teilbilder verwenden dieselbe Farbskala (0 bis zum größten Wert), sodass die
Farben direkt vergleichbar sind.

\begin{figure}[htbp]
  \centering
  \includegraphics[width=\textwidth]{fig4_spatial_heatmaps.png}
  \caption{Räumlicher Fehler des Standardmodells (L1-Verlust,
           \code{norm\_subject}, 200 Testbeispiele). Teilbilder 1--3:
           Gaußsche Fehler-Heatmaps ($r=0,15$, $\sigma=0,20$, gemeinsame
           Farbskala) in der Mitte $(0,5;\,0,5)$, am linken Rand
           $(0,1;\,0,5)$ und in der unteren linken Ecke $(0,1;\,0,1)$. Der
           Fehler wächst von der Mitte (0,086) über den Rand (0,425) zur Ecke
           (0,559). Teilbild 4: Ziele (blau) und Vorhersagen (rot); alle
           Vorhersagen liegen nahe der Bildschirmmitte.}
  \label{fig:heatmaps}
\end{figure}

Der Fehler nimmt von innen nach außen monoton zu: Der mittlere Fehler beträgt
$0,086$ in der Mittelregion ($n=18$), $0,432$ an den Rändern ($n=62$) und
$0,586$ in den Ecken ($n=42$). Die Randfehler sind etwa fünfmal und die
Eckfehler etwa siebenmal so groß wie die Fehler in der Mitte.

Das vierte Teilbild erklärt dieses Muster. Alle Vorhersagen liegen in einem
sehr kleinen Bereich um $(0,5;\,0,5)$: Ihre Standardabweichung beträgt in
beiden Koordinaten etwa $0,004$, während die Ziele eine Standardabweichung
von etwa $0,28$ haben. Das Modell sagt daher für jedes Bild fast denselben
Punkt vorher, und sein Fehler ist einfach der Abstand des Ziels zur
Bildschirmmitte. Ein konstanter Mittelpunkt-Schätzer ergibt praktisch dieselben
Werte (Mitte $0,085$, Ränder $0,432$, Ecken $0,587$; insgesamt
$\Ediag=\SI{28,45}{\percent}$ gegenüber \SI{28,47}{\percent} für das Modell).

\subsection{Verlustfunktionen}
\label{sec:losses}

Die hohen Rand- und Eckfehler führten zu der Idee, diesen Regionen im Verlust
mehr Gewicht zu geben. Sei $\hat{\bm p}_i=(\hat x_i,\hat y_i)$ die Vorhersage
und $\bm p_i=(x_i,y_i)$ das Ziel des Beispiels $i$ in einem Batch der Größe
$N$.

\paragraph{L1-Verlust (\code{nn.L1Loss}, \enquote{standard}).}
\begin{equation}
  \mathcal{L}_{1} = \frac{1}{2N}\sum_{i=1}^{N}\bigl(|\hat x_i-x_i| + |\hat y_i-y_i|\bigr).
\end{equation}
Der Gradient hat einen konstanten Betrag, unabhängig von der Fehlergröße. Der
Verlust ist daher robust gegenüber Ausreißern (zum Beispiel falschen Labels
während Zielübergängen), behandelt aber alle Bildschirmpositionen gleich.

\paragraph{Smooth-L1-Verlust (\code{nn.SmoothL1Loss}, $\beta=1$).}
Für ein Residuum $u$ gilt:
\begin{equation}
  \ell_\beta(u) =
  \begin{cases}
    \dfrac{u^2}{2\beta}, & |u| < \beta,\\[1.2ex]
    |u| - \dfrac{\beta}{2}, & \text{sonst.}
  \end{cases}
\end{equation}
Er ist quadratisch für kleine und linear für große Residuen. Da alle Residuen
in $[0,1]^2$ kleiner als $\beta=1$ sind, verhält er sich hier wie ein
skalierter quadratischer Fehler: Der Gradient wird klein, wenn der Fehler
klein ist.

\paragraph{Sensitiver L1-Verlust (positionsgewichtet, \code{sensitive}).}
Jedes Beispiel wird mit seinem $L_1$-Abstand zur Bildschirmmitte gewichtet:
\begin{equation}
  d_i = |x_i-0,5| + |y_i-0,5| \in [0,1],\qquad
  w_i = 1 + d_i \in [1,2],
\end{equation}
\begin{equation}
  \mathcal{L}_{\mathrm{sens}} = \frac{1}{N}\sum_{i=1}^{N}
  w_i\cdot\frac{|\hat x_i-x_i| + |\hat y_i-y_i|}{2}.
\end{equation}
Ein Eckbeispiel hat das doppelte Gewicht eines Beispiels in der Mitte, ein
Randbeispiel etwa das $1,4$-fache. Das vergrößert den Gradienten für äußere
Ziele und soll die Vorhersagen von der Mitte wegziehen. (Der frühere Entwurf
definierte das Gewicht als $w=1/d$; das entspricht nicht der Implementierung
in \code{model.py} und wäre in der Mitte unendlich.)

Der \enquote{verbesserte sensitive Verlust} in \cref{tab:loss_aux} war eine
weitere Variante aus einem explorativen Notebook. Seine genaue Formel ist im
aktuellen Code nicht enthalten; er wird daher nur der Vollständigkeit halber
angegeben.

\subsection{Ergebnisse des Verlustvergleichs}
\label{sec:loss_results}

\Cref{tab:loss} verbindet den globalen Vergleich aus der Variation mit fünf
Läufen mit den räumlichen Fehlern der in \cref{sec:partition} definierten
Regionen. Sie ersetzt die beiden getrennten Tabellen des früheren Entwurfs
(globale Verlusttabelle und $3\times3$-Rastertabelle).

\begin{table}[htbp]
  \centering
  \caption{Vergleich von L1- und sensitivem L1-Verlust. Links: mittlerer
           Testfehler $\Ediag$ über fünf Läufe (\code{exp.py},
           \SI{98}{\percent}-KI) und mittlere Laufzeit. Rechts: mittlerer
           euklidischer Fehler pro Bildschirmregion (normierte Einheiten; je
           ein früherer Lauf, 200 Testbeispiele, $r=0,15$; $n$ = Anzahl der
           Beispiele). Letzte Zeile: konstanter Mittelpunkt-Schätzer auf
           denselben 200 Testbeispielen.}
  \label{tab:loss}
  \small
  \setlength{\tabcolsep}{4.5pt}
  \begin{tabular}{lccc@{\hspace{1.2em}}ccc}
    \toprule
    & \multicolumn{3}{c}{Global (5 Läufe)} & \multicolumn{3}{c}{Räumlicher Fehler (1 Lauf)} \\
    \cmidrule(lr){2-4}\cmidrule(lr){5-7}
    Verlust & $\Ediag$ (\%) & \SI{98}{\percent}-KI (\%) & Laufzeit (min)
            & Mitte & Ränder & Ecken \\
            &               &                           &
            & ($n=18$) & ($n=62$) & ($n=42$) \\
    \midrule
    L1 (standard)\sel{}    & \textbf{28,32} & [28,23; 28,42] & 0,011 & 0,086 & 0,432 & 0,586 \\
    Sensitiver L1          & 28,48 & [28,19; 28,78] & 0,021 & 0,090 & 0,431 & 0,585 \\
    \midrule
    Konstante Mitte        & 28,45\textsuperscript{a} & --- & --- & 0,085 & 0,432 & 0,587 \\
    \bottomrule
  \end{tabular}\\[1mm]
  \parbox{0.95\linewidth}{\footnotesize \textsuperscript{a}\,Ausgewertet auf
  denselben 200 Testbeispielen wie die räumliche Analyse.}
\end{table}

\begin{table}[htbp]
  \centering
  \caption{Hilfsexperiment: räumlicher Fehler von vier Verlustfunktionen nach
           15 Epochen auf einer Trainingsmenge mit 52 Eckbeispielen (zwei
           Läufe, A und B, Werte wie protokolliert). Diese Zahlen stammen aus
           einem explorativen Notebook; die Berechnung der Regionswerte dieses
           Laufs ist im aktuellen Code nicht enthalten, daher sind sie mit
           \cref{tab:loss} nicht direkt vergleichbar.}
  \label{tab:loss_aux}
  \small
  \setlength{\tabcolsep}{5pt}
  \begin{tabular}{l cc cc cc cc}
    \toprule
    & \multicolumn{2}{c}{Gesamt (\%)} & \multicolumn{2}{c}{Mitte (\%)}
    & \multicolumn{2}{c}{Ränder (\%)} & \multicolumn{2}{c}{Ecken (\%)} \\
    \cmidrule(lr){2-3}\cmidrule(lr){4-5}\cmidrule(lr){6-7}\cmidrule(lr){8-9}
    Verlust & A & B & A & B & A & B & A & B \\
    \midrule
    L1 (\code{nn.L1Loss})             & 24,54 & 24,63 & 17,01 & 16,74 & 28,28 & 28,98 & 34,93 & 35,07 \\
    Smooth L1 (\code{nn.SmoothL1Loss})& 23,94 & 24,91 & 15,40 & 16,41 & 29,78 & 30,59 & 34,13 & 35,17 \\
    Sensitiver L1                     & 25,57 & 25,55 & 19,11 & 18,61 & 29,87 & 29,93 & 33,41 & 34,20 \\
    Verbesserter sensitiver L1        & 25,53 & 27,75 & 19,20 & 21,86 & 29,17 & 35,03 & 33,77 & 31,58 \\
    \bottomrule
  \end{tabular}
\end{table}

\paragraph{Beobachtungen.}
\begin{itemize}[nosep]
  \item Global ist L1 leicht, aber signifikant besser als der sensitive
        Verlust: \SI{28,32}{\percent} gegenüber \SI{28,48}{\percent}
        (Mann-Whitney $p=0,032$, \cref{tab:loss}); der L1-Wert ist daher fett
        gedruckt. Der sensitive Verlust streut außerdem stärker zwischen den
        Läufen. Beide Werte liegen nahe an der Konstant-Baseline
        (\SI{28,46}{\percent}).
  \item Räumlich verändert die Positionsgewichtung die Regionsfehler um
        weniger als $0,005$ normierte Einheiten. Die Rand- und Eckfehler
        bleiben etwa fünf- bis siebenmal so hoch wie der Fehler in der Mitte.
  \item Beide Modelle haben dasselbe räumliche Profil wie ein konstanter
        Mittelpunkt-Schätzer. Die Gewichtung vergrößert den Gradienten für
        äußere Beispiele, das reicht aber nicht, damit das Modell die
        Mittelpunktlösung verlässt.
  \item Im Hilfsexperiment (\cref{tab:loss_aux}) senken die sensitiven
        Varianten den Eckfehler leicht (sensitiver L1 um
        \num{0,9}--\num{1,5}, die verbesserte Variante um
        \num{1,2}--\num{3,5} Prozentpunkte), erhöhen aber den Fehler in der
        Mitte um \num{1,9}--\num{5,1} Prozentpunkte und den Gesamtfehler in
        allen Läufen. Die beiden Läufe unterscheiden sich deutlich, daher sind
        diese Tendenzen nicht verlässlich.
\end{itemize}

\paragraph{Fazit.}
Die räumliche Analyse hat gezeigt, \emph{was} das Problem ist: Die Fehler
konzentrieren sich an den Rändern und in den Ecken. Der positionsgewichtete
Verlust war der Versuch, diese Fehler zu verringern. Er hat die Ecken nicht
verbessert und war global sogar leicht, aber signifikant schlechter als L1
($p=0,032$); im Hilfsexperiment erhöhte er den Fehler in der Mitte. Der
Standard-L1-Verlust wird daher beibehalten. Er ergab den niedrigeren Fehler,
ist einfacher, hat keinen zusätzlichen Hyperparameter, und sein konstanter
Gradientenbetrag macht ihn robust gegenüber Label-Rauschen in den
persönlichen Aufnahmen. Die Analyse zeigt auch, \emph{warum} die Gewichtung
nicht helfen konnte: Die großen äußeren Fehler entstehen nicht durch eine
falsche Gewichtung der Bildschirmregionen, sondern dadurch, dass das Modell
die Bildschirmmitte kaum verlässt (\cref{sec:discussion}).

% ==============================================================================
\section{Personalisierungssitzungen}
\label{sec:sessions}
% ==============================================================================

\subsection{Aufbau}
\label{sec:sessions_setup}

In jeder persönlichen Sitzung wird der Kopf eines GazeCapture-Basismodells
auf den Frames einer Aufnahme neu trainiert (nur FC, Sigmoid, L1,
$\eta=10^{-4}$, Weight Decay $10^{-5}$, Batchgröße 32, AdamW). Sofern nicht
anders angegeben, wird jede Sitzung zufällig in \SI{70}{\percent} Trainings-,
\SI{15}{\percent} Validierungs- und \SI{15}{\percent} Testframes aufgeteilt
(Seed 42). Das Early Stopping (Patience 3) verwendet den Validierungsfehler,
und angegeben wird der Testfehler in der Epoche mit dem niedrigsten
Validierungsfehler. Ist die beste Epoche~0, hat die Personalisierung den
Validierungsfehler nicht verbessert, und der angegebene Testfehler ist der des
unveränderten Basismodells.

Alle Ergebnisse dieses und der folgenden Abschnitte wurden mit dem Skript
\code{personal\_exp.py} (Notebook \code{personal\_compare.ipynb}) und der
korrigierten Metrik (\cref{eq:ediag}) erzeugt. Alle Zufalls-Seeds sind fest,
und der Backbone bleibt einschließlich seiner BatchNorm-Statistiken
vollständig eingefroren.

\paragraph{Merkmals-Cache.} Da der Backbone eingefroren ist, bildet er ein
Bild in jeder Epoche auf dieselben 512 Merkmale ab. Diese Merkmale werden
einmal pro Sitzung und Basismodell berechnet (Evaluationstransformation, ohne
Augmentierung) und gespeichert; danach wird nur noch der Regressionskopf
trainiert. Das verändert das trainierte Modell nicht, sorgt aber dafür, dass
ein Lauf auf einer CPU nur Sekunden dauert. Der einzige Unterschied zu den
früheren Notebook-Läufen ist, dass während der Personalisierung keine
Bildaugmentierung angewendet wird.

\paragraph{Basismodelle (E0).} Die beiden Basismodelle werden mit der in
\cref{sec:training} gewählten Konfiguration (nur FC, Sigmoid, L1,
$\eta=10^{-4}$, Batchgröße 32, \num{1000}/\num{200} zufällige Frames, Seed~0,
validierungsbasiertes Early Stopping) auf \code{norm\_subject} und
\code{norm\_random} trainiert. Von fünf Seeds wird der Lauf mit dem
niedrigsten \emph{Validierungs}fehler behalten (\code{train\_base.py},
\code{models/base/base\_<split>.pt} mit einer JSON-Datei, die die
Konfiguration und alle Seeds dokumentiert). Für \code{norm\_subject} wurde
Seed~0 gewählt (Validierung \SI{27,37}{\percent}, Test \SI{28,70}{\percent};
Konstant-Baseline derselben Testdaten \SI{28,46}{\percent}), für
\code{norm\_random} Seed~3 (Validierung \SI{26,74}{\percent}, Test
\SI{27,81}{\percent}; Konstant-Baseline \SI{27,66}{\percent}). Die fünf Seeds
unterscheiden sich um weniger als \num{0,4} Prozentpunkte, die Wahl des Seeds
ist also unkritisch. Wie in \cref{sec:discussion} liegen beide Basismodelle
auf dem Niveau der Konstant-Baseline.

Von den 27 aufgenommenen Sitzungen werden 18 verwendet: Gen~1 (01--04), Gen~2
(05--09) und Gen~4 (19--27). Gen~3 (10--18) wurde verworfen, weil sie keinen
Lerneffekt zeigte. \Cref{tab:session_conditions} listet die
Aufnahmebedingungen auf.

\begin{table}[htbp]
  \centering
  \caption{Aufnahmebedingungen der 18 Sitzungen. Gen~1: kontinuierliche
           Zielbewegung ohne Sprünge; Gen~2: Ziel mit zufälligen Sprüngen;
           Gen~4: 5-s-Segmente (\SI{0,5}{\second} WAIT an einem neuen
           Startpunkt, \SI{4,5}{\second} MOVE), \SI{90}{\percent} der
           Startpunkte in den Bildschirmecken. Dauer und Betriebssystem wurden
           nur in Gen~4 variiert (Gen~1/2: Ubuntu, Dauer nicht protokolliert).}
  \label{tab:session_conditions}
  \small
  \begin{tabular}{llll}
    \toprule
    Sitzungen & Generation & Brille & Licht / Hintergrund \\
    \midrule
    01, 03, 04 & Gen 1 & mit    & 01, 04: hell; 03: dunkel \\
    02         & Gen 1 & ohne   & hell \\
    05, 07, 08 & Gen 2 & mit    & 05: hell; 07: dunkel; 08: sehr dunkel \\
    06, 09     & Gen 2 & ohne   & 06: hell; 09: sehr dunkel \\
    \midrule
    \multicolumn{4}{l}{Gen 4 \quad (Dauer / BS)} \\
    19, 20 & Gen 4 & mit    & ruhig; 120\,s / 90\,s; Ubuntu \\
    21, 22 & Gen 4 & ohne   & ruhig; 120\,s / 90\,s; Ubuntu \\
    23 / 24 & Gen 4 & mit / ohne & unruhig; 120\,s; Ubuntu \\
    25 / 26 & Gen 4 & mit / ohne & unruhig; 120\,s; Windows \\
    27     & Gen 4 & mit    & unruhig; 90\,s; Windows \\
    \bottomrule
  \end{tabular}
\end{table}

\subsection{Ergebnisse aller Sitzungen (E1)}
\label{sec:sessions_results}

\Cref{tab:sessions} zeigt jede Sitzung mit beiden Basismodellen.

\begin{table}[htbp]
  \centering
  \caption{Personalisierung der 18 Sitzungen mit zwei Basismodellen
           (ursprüngliche Labels, zufälliger 70/15/15-Split, Seed 42). Basis:
           Testfehler des unveränderten Basismodells (Epoche~0); Val.:
           bester Validierungsfehler; Test: Testfehler in der besten
           Validierungsepoche; Ep.: beste Epoche; $t$: Dauer eines Laufs in
           Sekunden (\cref{sec:session_time}); Konst.: Konstant-Baseline
           (Mittelwert der Trainingsziele, ausgewertet auf denselben
           Testframes). Fehler in \si{\percent} (\cref{eq:ediag}).
           Gruppenzeilen: Mittelwert und \SI{98}{\percent}-KI über die
           Sitzungen der Gruppe.}
  \label{tab:sessions}
  \footnotesize
  \setlength{\tabcolsep}{2.6pt}
  \resizebox{\textwidth}{!}{%
  \begin{tabular}{cc ccccc ccccc c}
    \toprule
    & & \multicolumn{5}{c}{Basismodell \code{norm\_subject}}
      & \multicolumn{5}{c}{Basismodell \code{norm\_random}} & \\
    \cmidrule(lr){3-7}\cmidrule(lr){8-12}
    Sitzung & Gen. & Basis & Val. & Test & Ep. & $t$ (s) & Basis & Val. & Test & Ep. & $t$ (s) & Konst. \\
    \midrule
    01 & 1 & 29,19 & 14,82 & 15,83 & 22 & 2,3 & 29,49 & 11,70 & 12,97 & 47 & 4,2 & 29,45 \\
    02 & 1 & 28,48 & 10,08 & 10,48 & 30 & 2,2 & 28,42 & 10,04 & 10,43 & 44 & 3,2 & 28,49 \\
    03 & 1 & 29,61 & 15,26 & 17,20 & 33 & 2,5 & 30,13 & 15,41 & 17,35 & 36 & 2,6 & 29,76 \\
    04 & 1 & 29,98 & 16,88 & 17,64 & 17 & 2,1 & 29,38 & 15,23 & 16,16 & 31 & 2,7 & 29,17 \\
    05 & 2 & 32,92 & 25,43 & 17,64 & 26 & 1,1 & 32,49 & 31,93 & 32,77 & 2 & 0,2 & 32,72 \\
    06 & 2 & 29,29 & 11,31 & 16,67 & 54 & 1,8 & 28,89 & 12,09 & 16,55 & 55 & 2,0 & 27,02 \\
    07 & 2 & 32,18 & 24,84 & 30,99 & 16 & 0,7 & 32,43 & 25,02 & 31,62 & 18 & 0,6 & 32,28 \\
    08 & 2 & 25,82 & 23,65 & 23,30 & 24 & 0,8 & 26,32 & 23,38 & 22,31 & 23 & 0,6 & 27,00 \\
    09 & 2 & 26,83 & 23,55 & 28,64 & 4 & 0,2 & 26,21 & 15,19 & 16,87 & 47 & 1,9 & 28,17 \\
    19 & 4 & 30,42 & 13,38 & 17,76 & 33 & 2,8 & 30,04 & 12,13 & 16,85 & 47 & 4,4 & 30,19 \\
    20 & 4 & 33,14 & 16,06 & 14,64 & 23 & 1,9 & 33,01 & 14,58 & 12,93 & 35 & 2,7 & 33,49 \\
    21 & 4 & 30,18 & 15,03 & 12,51 & 39 & 4,6 & 30,46 & 15,80 & 13,20 & 23 & 2,6 & 29,97 \\
    22 & 4 & 31,17 & 13,54 & 13,02 & 22 & 2,2 & 30,87 & 12,53 & 12,32 & 37 & 3,1 & 30,32 \\
    23 & 4 & 33,15 & 23,48 & 23,65 & 10 & 1,4 & 32,59 & 24,02 & 24,05 & 10 & 2,0 & 32,57 \\
    24 & 4 & 33,09 & 28,77 & 26,97 & 7 & 0,9 & 32,84 & 29,53 & 27,42 & 6 & 1,0 & 33,01 \\
    25 & 4 & 32,19 & 26,98 & 24,61 & 19 & 2,4 & 33,00 & 26,74 & 24,13 & 26 & 2,8 & 31,54 \\
    26 & 4 & 30,60 & 21,93 & 18,61 & 20 & 2,3 & 30,39 & 21,40 & 17,97 & 37 & 4,5 & 29,48 \\
    27 & 4 & 31,98 & 26,05 & 26,06 & 19 & 1,5 & 32,29 & 26,27 & 26,68 & 19 & 1,7 & 31,19 \\
    \midrule
    \multicolumn{2}{l}{Gen 1 (4)} & 29,32 & 14,26 & 15,29 & 25,5 & 2,3 & 29,35 & 13,10 & 14,23 & 39,5 & 3,2 & 29,22 \\
      & & [27,86; 30,77] & [7,62; 20,89] & [7,80; 22,77] & & [1,8; 2,7] & [27,75; 30,96] & [7,07; 19,12] & [7,10; 21,35] & & [1,5; 4,8] & [27,99; 30,45] \\
    \multicolumn{2}{l}{Gen 2 (5)} & 29,41 & 21,76 & 23,45 & 24,8 & 0,9 & 29,27 & 21,52 & 24,02 & 29,0 & 1,0 & 29,44 \\
      & & [24,14; 34,68] & [11,88; 31,63] & [12,74; 34,16] & & [0,0; 1,9]\textsuperscript{a} & [24,06; 34,47] & [8,20; 34,85] & [10,93; 37,12] & & [0,0; 2,5]\textsuperscript{a} & [24,68; 34,20] \\
    \multicolumn{2}{l}{Gen 4 (9)} & 31,77 & 20,58 & 19,76 & 21,3 & 2,2 & 31,72 & 20,33 & 19,51 & 26,7 & 2,8 & 31,31 \\
      & & [30,60; 32,94] & [14,66; 26,50] & [14,26; 25,26] & & [1,2; 3,3] & [30,51; 32,93] & [13,88; 26,79] & [13,59; 25,42] & & [1,7; 3,9] & [29,91; 32,70] \\
    \midrule
    \multicolumn{2}{l}{\textbf{Alle (18)}} & \textbf{30,57} & \textbf{19,50} & \textbf{19,79} & \textbf{23,2} & \textbf{1,9} & \textbf{30,51} & \textbf{19,05} & \textbf{19,59} & \textbf{30,2} & \textbf{2,4} & \textbf{30,32} \\
      & & [29,27; 31,87] & [15,87; 23,13] & [16,18; 23,40] & & [1,3; 2,5] & [29,21; 31,82] & [14,86; 23,25] & [15,47; 23,71] & & [1,6; 3,2] & [29,12; 31,53] \\
    \bottomrule
  \end{tabular}}\\[1mm]
  \parbox{0.95\linewidth}{\footnotesize \textsuperscript{a}\,Untere Grenze
  des $t$-Intervalls unter null (wenige Sitzungen, große Streuung), bei 0
  abgeschnitten.}
\end{table}

\begin{figure}[htbp]
  \centering
  \includegraphics[width=\textwidth]{fig_sessions_baseline.png}
  \caption{Testfehler der 18 Sitzungen nach der Personalisierung mit den
           beiden Basismodellen im Vergleich mit der Konstant-Baseline und mit
           dem unveränderten Basismodell (\code{norm\_subject}) auf denselben
           Testframes (Daten aus \cref{tab:sessions}).}
  \label{fig:sessions}
\end{figure}

\paragraph{Die Personalisierung wirkt.}
Über alle 18 Sitzungen sinkt der Testfehler von \SI{30,32}{\percent}
(Konstant-Baseline) auf \SI{19,79}{\percent} mit dem Basismodell
\code{norm\_subject} und auf \SI{19,59}{\percent} mit dem Basismodell
\code{norm\_random}, eine Verringerung um \SI{35}{\percent} bzw.\
\SI{36}{\percent} (Wilcoxon-Vorzeichen-Rang-Test $p=1,5\cdot10^{-5}$ für
beide; gepaarter $t$-Test $p<10^{-6}$). Die unveränderten Basismodelle sind
nicht besser als die Konstant-Baseline (\SI{30,57}{\percent} und
\SI{30,51}{\percent}; Wilcoxon $p=0,25$ und $p=0,35$), der gesamte Gewinn
stammt also aus den wenigen Minuten persönlicher Daten. Die beste Sitzung,
S02 (Gen~1, heller Raum, ohne Brille), erreicht \SI{10,5}{\percent}. Die
neuen Läufe bestätigen die früheren Notebook-Läufe (Mittelwert
\SI{19,85}{\percent} über dieselben Sitzungen mit dem Basismodell
\code{norm\_subject}).

\paragraph{Gruppen.} Die Verringerung gegenüber der Konstant-Baseline beträgt
\SI{48}{\percent} / \SI{51}{\percent} in Gen~1, \SI{20}{\percent} /
\SI{18}{\percent} in Gen~2 und \SI{37}{\percent} / \SI{38}{\percent} in Gen~4
(\code{norm\_subject} / \code{norm\_random}). In Gen~4 ist die Verbesserung
signifikant (Wilcoxon $p=0,004$); für Gen~1 ($n=4$) und Gen~2 ($n=5$) beträgt
der kleinstmögliche $p$-Wert des Tests $0,125$ bzw.\ $0,0625$, sodass diese
Gruppen für ein signifikantes Ergebnis zu klein sind, obwohl der Effekt in
Gen~1 groß ist. Die Unterschiede zwischen den Generationen sind nicht
signifikant (Kruskal-Wallis $p=0,13$ und $p=0,12$) und dürfen nicht allein
als Wirkung des Aufnahmeprotokolls gelesen werden
(\cref{sec:session_design}).

\paragraph{Wenn die Personalisierung scheitert.}
In zwei Läufen ist das personalisierte Modell nicht besser als die
Konstant-Baseline: S09 mit \code{norm\_subject} (beste Epoche~4,
\SI{28,64}{\percent} gegenüber \SI{28,17}{\percent}) und S05 mit
\code{norm\_random} (beste Epoche~2, \SI{32,77}{\percent} gegenüber
\SI{32,72}{\percent}). In beiden Fällen hörte der Validierungsfehler sehr früh
auf, sich zu verbessern. Mit dem jeweils anderen Basismodell erreichen
dieselben Sitzungen \SI{16,87}{\percent} (S09) und \SI{17,64}{\percent} (S05).
Das Scheitern ist daher keine Eigenschaft der Sitzung, sondern eines
einzelnen Laufs: Mit nur 20--30 Validierungsframes (Gen~2) kann das Early
Stopping das Training nach wenigen Epochen beenden. S07 (dunkler Raum)
verbessert sich mit beiden Basismodellen nur wenig (\SI{31,0}{\percent} /
\SI{31,6}{\percent} gegenüber \SI{32,3}{\percent}); gute Beleuchtung bleibt
eine Voraussetzung für die Webcam-Personalisierung.

\paragraph{Basismodell.}
Beide Basismodelle ergeben dasselbe Ergebnis (\SI{19,79}{\percent} gegenüber
\SI{19,59}{\percent}; Wilcoxon $p=0,20$). Das passt zum Hauptergebnis, dass
das Basismodell vor allem eine Initialisierung liefert und die persönlichen
Daten das Ergebnis bestimmen. Für die folgenden Experimente wird das
Basismodell \code{norm\_subject} verwendet, weil es ohne Informationen über
die Testpersonen trainiert wurde.

\paragraph{Validierungs- und Testfehler.} Im Mittel liegt der
Validierungsfehler (\SI{19,50}{\percent}) nahe am Testfehler
(\SI{19,79}{\percent}). Die größten Abweichungen (5--8 Prozentpunkte,
S05--S07, S09) treten in den kurzen Gen-2-Aufnahmen auf, in denen die
Validierungs- und Testmengen jeweils nur etwa 25 Frames enthalten.

\subsection{Trainingszeit einer Sitzung}
\label{sec:session_time}

Die Zeit $t$ in \cref{tab:sessions} ist die tatsächlich verstrichene Dauer
\emph{eines vollständigen Personalisierungslaufs} einer Sitzung auf den
zwischengespeicherten Merkmalen: die Auswertung des unveränderten Modells
(Epoche~0), alle Trainingsepochen bis zum Early Stopping und die Auswertung
der Trainings-, Validierungs- und Testdaten nach jeder Epoche. Sie ist weder
eine Zeit pro Epoche noch ein Mittelwert über mehrere Läufe. Sie enthält
nicht die einmalige Extraktion der Backbone-Merkmale der Sitzung, die alle
Läufe derselben Sitzung und desselben Basismodells gemeinsam nutzen. Ein Lauf
dauert auf einer CPU im Mittel \SI{1,9}{\second} (\code{norm\_subject}) bzw.\
\SI{2,4}{\second} (\code{norm\_random}); die Dauer wächst mit der Zahl der
Frames und der Zahl der Epochen (kurze Gen-2-Aufnahmen: \SI{0,9}{\second}).
Die früheren Notebook-Läufe, die jedes Bild in jeder Epoche durch den
Backbone schickten, benötigten 1--26 Minuten pro Sitzung. Die
\enquote{Laufzeit} der Hyperparametertabellen (\cref{sec:protocol}) ist
dagegen der Mittelwert über fünf wiederholte Läufe einer
GazeCapture-Konfiguration.

\subsection{Label-Bereinigung (Gen 4, E1b)}
\label{sec:cleaning}

In Gen~4 beginnt jedes Segment mit einer WAIT-Phase von \SI{0,5}{\second}, in
der das Ziel an einer neuen Position erscheint und der Nutzer es erst finden
muss. Diese Frames (\SI{10,0}{}--\SI{10,5}{\percent} der Labels) können ein
falsches Blicklabel haben. Sie sind exakt protokolliert, daher lassen sich
bereinigte Labeldateien ohne die WAIT-Frames schreiben
(\code{session\_design.py --write-clean}).

Ein einfacher Vergleich ist nicht fair: Das bereinigte Modell wird ohne
WAIT-Frames getestet, das ursprüngliche Modell dagegen auf allen Frames,
einschließlich der schwer vorherzusagenden WAIT-Frames. Deshalb wird das
ursprüngliche Modell zusätzlich nur auf den MOVE-Frames seines Testsplits
ausgewertet (dasselbe Modell, WAIT-Frames aus der Testmenge entfernt;
\cref{tab:cleaning}).

\begin{table}[htbp]
  \centering
  \caption{Label-Bereinigung in Gen~4 (Basismodell \code{norm\_subject},
           zufälliger Split, Seed 42). Ursprünglich: auf allen Labels
           trainiert; Testfehler auf allen Testframes, nur auf den MOVE-Frames
           und nur auf den WAIT-Frames ($n_W$ = Anzahl der WAIT-Testframes).
           Bereinigt: ohne WAIT-Frames trainiert und getestet. $\Delta$ =
           Bereinigt $-$ Ursprünglich (MOVE).}
  \label{tab:cleaning}
  \footnotesize
  \setlength{\tabcolsep}{4pt}
  \begin{tabular}{c cccc cc c}
    \toprule
    & \multicolumn{4}{c}{Ursprüngliche Labels} & \multicolumn{2}{c}{Bereinigte Labels} & \\
    \cmidrule(lr){2-5}\cmidrule(lr){6-7}
    Sitzung & Alle & MOVE & WAIT & $n_W$ & Test & Konst. & $\Delta$ \\
    \midrule
    19 & 17,76 & 13,28 & 41,12 & 8  & 13,06 & 29,07 & $-0,22$ \\
    20 & 14,64 & 12,24 & 28,38 & 8  & 12,77 & 30,71 & $+0,53$ \\
    21 & 12,51 &  9,77 & 27,59 & 10 &  9,79 & 29,23 & $+0,01$ \\
    22 & 13,02 & 10,21 & 27,64 & 8  &  9,96 & 31,71 & $-0,25$ \\
    23 & 23,65 & 20,95 & 41,84 & 10 & 17,18 & 30,92 & $-3,77$ \\
    24 & 26,97 & 24,96 & 41,93 & 10 & 23,99 & 31,64 & $-0,97$ \\
    25 & 24,61 & 21,63 & 44,36 & 10 & 17,94 & 30,49 & $-3,70$ \\
    26 & 18,61 & 14,34 & 41,72 & 10 & 14,58 & 30,20 & $+0,24$ \\
    27 & 26,06 & 23,10 & 44,80 & 8  & 23,69 & 30,59 & $+0,60$ \\
    \midrule
    Mittel & 19,76 & 16,72 & 37,71 & 9,1 & 15,88 & 30,51 & $-0,84$ \\
    \bottomrule
  \end{tabular}
\end{table}

\paragraph{Ergebnis.} Die WAIT-Frames werden sehr schlecht vorhergesagt
(\SI{37,7}{\percent}, über der Konstant-Baseline): Der Blick ist noch nicht
auf dem Ziel, ihre Labels sind also falsch. Werden sie nur aus der
\emph{Test}menge entfernt, sinkt der Fehler von \SI{19,76}{\percent} auf
\SI{16,72}{\percent}. Werden sie auch aus den \emph{Trainings}daten entfernt,
sinkt er weiter auf \SI{15,88}{\percent}, ein mittlerer Gewinn von \num{0,84}
Prozentpunkten, der nicht signifikant ist (Wilcoxon $p=0,43$). Der Gewinn
stammt aus zwei Sitzungen mit unruhigem Hintergrund (S23, S25, etwa $-3,7$
Punkte); in den ruhigen Sitzungen liegt der Unterschied unter \num{0,6}
Punkten. Fazit: Die WAIT-Frames sollten zumindest aus der Auswertung entfernt
werden, weil sie den Labelfehler und nicht den Modellfehler messen; sie auch
aus dem Training zu entfernen, kostet nichts und hilft in schwierigen
Sitzungen, der Effekt ist aber klein.

\paragraph{Variabilität von Lauf zu Lauf.} Ein früherer, ungültiger Vergleich
bestand aus zwei Notebook-Läufen derselben Konfiguration auf identischen Daten
(Gen~4, Zufalls-Seeds nicht fest). Sie unterschieden sich im Mittel um
\num{1,9} und höchstens um \num{5,5} Prozentpunkte. Die neuen Läufe mit drei
festen Seeds (\cref{tab:loading}) zeigen eine Standardabweichung von bis zu
\num{3,1} Prozentpunkten zwischen den Seeds. Unterschiede zwischen
\emph{einzelnen} Läufen in dieser Größenordnung sind daher nicht
interpretierbar; das erklärt auch die oben genannten Fälle S05/S09.

\subsection{Aufnahmebedingungen (Gen 4)}
\label{sec:conditions}

Gen~4 erlaubt Vergleiche, in denen sich nur ein Faktor ändert. Jeder Wert ist
die Differenz zweier einzelner Sitzungen (Testfehler aus \cref{tab:sessions},
\code{norm\_subject} / \code{norm\_random}, Prozentpunkte; positiv = die erste
Bedingung ist schlechter):
\begin{itemize}[nosep]
  \item \textbf{Brille} (mit $-$ ohne): $+5,2/+3,6$ (S19--S21),
        $+1,6/+0,6$ (S20--S22), $-3,3/-3,4$ (S23--S24), $+6,0/+6,2$ (S25--S26);
  \item \textbf{Dauer} (\SI{90}{\second} $-$ \SI{120}{\second}):
        $-3,1/-3,9$ (S20--S19), $+0,5/-0,9$ (S22--S21);
  \item \textbf{Hintergrund} (unruhig $-$ ruhig): $+5,9/+7,2$ (S23--S19),
        $+14,5/+14,2$ (S24--S21);
  \item \textbf{Betriebssystem} (Windows $-$ Ubuntu): $+1,0/+0,1$
        (S25--S23), $-8,4/-9,4$ (S26--S24).
\end{itemize}
Der einzige einheitliche und große Effekt ist der \textbf{Hintergrund}: Ein
unruhiger Hintergrund erhöht den Fehler mit beiden Basismodellen um 6--15
Prozentpunkte. Eine Brille erhöht den Fehler tendenziell (drei von vier
Paaren), die Aufnahmelänge (90 vs.\ 120\,s) zeigt keinen klaren Effekt, und
der Vergleich der Betriebssysteme ist widersprüchlich. Jede Bedingung ist
durch einzelne Sitzungen vertreten, daher sind diese Werte beschreibend und
nicht statistisch getestet.

% ==============================================================================
\section{Strategien zum Laden der Daten (E2)}
\label{sec:loading}
% ==============================================================================

Der zufällige Split aus \cref{sec:sessions} hat eine Schwäche: Benachbarte
Frames eines Videos sind fast identisch (6 Labels pro Sekunde), sodass die
Testmenge Beinahe-Duplikate von Trainingsframes enthält. Das ist eine Form
von Datenleck (Data Leakage), die den zufälligen Split optimistisch macht.
Der folgende Vergleich misst, wie groß dieser Effekt ist.

\paragraph{Sitzungen.} Verwendet werden die Sitzungen 21 und 22. Sie werden
nach einer festen Regel ausgewählt: Sie sind die einzigen zwei Sitzungen, die
sich nur in einem Faktor unterscheiden (Dauer, \SI{120}{\second} vs.\
\SI{90}{\second}), und wurden ansonsten unter den besten Bedingungen
aufgenommen (Gen~4, ohne Brille, ruhiger Hintergrund, Ubuntu, gleicher Tag).
Das macht sie für den Transfer zwischen zwei Videos geeignet, und beide haben
genug Frames (720 und 540) für zeitliche Blöcke.

\paragraph{Strategien.}
\begin{description}[nosep]
  \item[Zufällige Auswahl] (\code{--split random}): \SI{70}{\percent} /
        \SI{15}{\percent} / \SI{15}{\percent} der Frames werden zufällig aus
        der Sitzung gezogen (Referenz; enthält Beinahe-Duplikate).
  \item[Sequenzielle Auswahl] (\code{--split sequential}): Die Frames bleiben
        in Aufnahmereihenfolge; die ersten \SI{70}{\percent} werden für das
        Training, die nächsten \SI{15}{\percent} für die Validierung und die
        letzten \SI{15}{\percent} für den Test verwendet. Eine Sekunde (6
        Labels) am Anfang des Validierungs- und des Testblocks wird verworfen,
        damit keine benachbarten Frames auf beiden Seiten einer Blockgrenze
        liegen. Das entspricht dem realen Anwendungsfall: Der Nutzer
        kalibriert zuerst und verwendet den Tracker danach.
  \item[Training mit zwei Videos] (\code{--split two\_videos}):
        \SI{80}{\percent} der Frames der ersten Sitzung werden für das
        Training und die übrigen \SI{20}{\percent} (zufällig) für die
        Validierung verwendet; die zweite Sitzung dient \emph{nur} dem Test
        (alle Frames). Beide Richtungen werden ausgeführt (21$\to$22 und
        22$\to$21). Das misst, ob sich eine Kalibrierung auf eine neue
        Aufnahme derselben Person übertragen lässt.
\end{description}

\paragraph{Fairer Vergleich.} Alle Strategien verwenden dasselbe Basismodell
(\code{norm\_subject}), dieselben Trainingseinstellungen, Early Stopping auf
ihren eigenen Validierungsdaten und dieselben drei Seeds (42, 43, 44). Für
jeden Lauf werden die Konstant-Baseline (Mittelwert der
\emph{Trainings}ziele, ausgewertet auf den Testframes dieses Laufs) und der
Fehler des unveränderten Basismodells auf denselben Testframes angegeben,
damit die unterschiedlichen Testmengen vergleichbar sind. Kein Testframe wird
für das Training oder die Modellauswahl verwendet.

\begin{table}[htbp]
  \centering
  \caption{Strategien zum Laden der Daten (Sitzungen 21 und 22, Basismodell
           \code{norm\_subject}, ursprüngliche Labels). Mittelwert über drei
           Seeds; SD: Standardabweichung des Testfehlers über die Seeds;
           Basis: unverändertes Basismodell; Red.: Verringerung des
           Testfehlers relativ zur Konstant-Baseline; $n$: Trainings- /
           Validierungs- / Testframes; $t$: Dauer eines Laufs (s).}
  \label{tab:loading}
  \footnotesize
  \setlength{\tabcolsep}{3.5pt}
  \resizebox{\textwidth}{!}{%
  \begin{tabular}{l c cc c cc c c c}
    \toprule
    Strategie & Training $\to$ Test & Test (\%) & SD & Val.\ (\%) & Basis (\%) & Konst.\ (\%) & Red.\ (\%) & $n$ & $t$ (s) \\
    \midrule
    Zufällig    & 21 $\to$ 21 & 14,40 & 3,13 & 14,15 & 31,52 & 31,25 & 53,9 & 503 / 108 / 109 & 3,9 \\
                & 22 $\to$ 22 & 13,30 & 0,38 & 15,98 & 30,97 & 30,68 & 56,7 & 378 / 81 / 81   & 3,1 \\
    Sequenziell & 21 $\to$ 21 & 14,47 & 1,22 & 19,09 & 27,62 & 27,43 & 47,2 & 503 / 102 / 103 & 2,0 \\
                & 22 $\to$ 22 & 16,63 & 0,14 & 14,41 & 28,83 & 29,91 & 44,4 & 378 / 75 / 75   & 2,4 \\
    Zwei Videos & 21 $\to$ 22 & 16,79 & 0,13 & 15,67 & 31,05 & 32,17 & 47,8 & 576 / 144 / 540 & 5,0 \\
                & 22 $\to$ 21 & 17,70 & 0,42 & 14,33 & 32,03 & 33,13 & 46,6 & 432 / 108 / 720 & 3,2 \\
    \midrule
    \multicolumn{2}{l}{Zufällig (Mittel)}    & 13,85 & & & 31,25 & 30,97 & 55,3 & & \\
    \multicolumn{2}{l}{Sequenziell (Mittel)} & 15,55 & & & 28,22 & 28,67 & 45,8 & & \\
    \multicolumn{2}{l}{Zwei Videos (Mittel)} & 17,24 & & & 31,54 & 32,65 & 47,2 & & \\
    \bottomrule
  \end{tabular}}
\end{table}

\paragraph{Ergebnis.} Die erwartete Reihenfolge wird bestätigt: zufällig
(\SI{13,85}{\percent}) $<$ sequenziell (\SI{15,55}{\percent}) $<$ zwei Videos
(\SI{17,24}{\percent}). Mit 6 Läufen pro Strategie (zwei Sitzungen $\times$
drei Seeds) ist nur der Unterschied zwischen sequenziellem Training und
Training mit zwei Videos signifikant (Mann-Whitney U, $p=0,009$); zufällig
vs.\ sequenziell ($p=0,09$) und zufällig vs.\ zwei Videos ($p=0,06$) sind es
nicht. Da sich die Konstant-Baselines zwischen den Testmengen unterscheiden,
ist die relative Verringerung das fairere Maß: \SI{55}{\percent} beim
zufälligen Split, aber nur \SI{46}{\percent}--\SI{47}{\percent} bei den
realistischen Strategien. Der zufällige Split überschätzt den Nutzen der
Personalisierung also um etwa 8--9 Prozentpunkte der Verringerung.

\paragraph{Interpretation.}
\begin{itemize}[nosep]
  \item Der zufällige Split ist optimistisch, auf diesen beiden guten
        Sitzungen aber nur mäßig (1,7 Prozentpunkte). Bei schwierigen
        Sitzungen ist der Effekt viel größer (\cref{sec:design_performance}).
  \item Eine Kalibrierung überträgt sich auf eine zweite Aufnahme derselben
        Person unter denselben Bedingungen: Das Modell mit zwei Videos
        erreicht auf einem nie gesehenen Video
        \SI{17}{\percent}--\SI{18}{\percent} bei \SI{47}{\percent}
        Verringerung. Der Transfer kostet gegenüber dem zufälligen Split etwa
        3 Prozentpunkte.
  \item Die Ergebnisse mit zwei Videos schwanken am wenigsten zwischen den
        Seeds (SD unter 0,5), da die Testmenge groß ist (540--720 Frames).
        Der zufällige Split hat die größte Streuung (SD 3,1 in S21), weil
        seine 81--109 Testframes vom Seed abhängen.
  \item Der Validierungsfehler des sequenziellen Splits sagt seinen
        Testfehler schlecht voraus (S21: \SI{19,1}{\percent} Validierung vs.\
        \SI{14,5}{\percent} Test), weil Validierungs- und Testblock
        unterschiedliche Bildschirmbereiche abdecken.
\end{itemize}
Für den Anwendungsfall eines Eye-Trackers, der einmal kalibriert und danach
benutzt wird, sind die Werte des sequenziellen Splits bzw.\ der zwei Videos
die realistische Schätzung.

% ==============================================================================
\section{Gestaltung der Personalisierungssitzung}
\label{sec:session_design}
% ==============================================================================

Die Gestaltung der Sitzung bestimmt, welche Bildschirmpositionen durch die
Kalibrierdaten abgedeckt werden. Da die Fehler an den Rändern und in den Ecken
am größten sind (\cref{sec:spatial}), soll die Gestaltung diese Bereiche gut
abdecken. Es wurden drei Varianten aufgenommen:
\begin{itemize}[nosep]
  \item \textbf{Standardvideo} (Gen~1, \code{base\_Code.py}): Das Ziel bewegt
        sich kontinuierlich und ohne Sprünge über den Bildschirm;
  \item \textbf{Zufällige Sprungpunkte} (Gen~2): Das Ziel springt alle
        \SI{1,4}{}--\SI{2,7}{\second} zu einem gleichverteilt gezogenen neuen
        Startpunkt;
  \item \textbf{Eckengewichtete Sprungpunkte} (Gen~4,
        \code{base\_Code\_rand.py}): 5-s-Segmente mit \SI{0,5}{\second} WAIT
        an einem neuen Startpunkt und \SI{4,5}{\second} MOVE;
        \SI{90}{\percent} der Startpunkte liegen in den vier Eckbereichen (je
        \SI{10}{\percent}\,$\times$\,\SI{10}{\percent} des Bildschirms).
\end{itemize}

\subsection{Räumliche Abdeckung und Übergangsframes}

\Cref{tab:design} wird direkt aus den Label- und Positionsdateien der 18
Sitzungen berechnet (\code{session\_design.py}); ein Training ist nicht
nötig. Die Regionsanteile verwenden die Bildschirmeinteilung aus
\cref{sec:partition} ($r=0,15$). Übergangsframes sind Labels, die
aufgenommen wurden, während sich der Blick noch zu einem neuen Ziel bewegt:
in Gen~4 die protokollierte WAIT-Phase, in Gen~2 die Labels innerhalb von
\SI{0,5}{\second} nach einem Sprung (geschätzt, keine Phasenprotokollierung),
in Gen~1 keine.

\begin{table}[htbp]
  \centering
  \caption{Räumliche Abdeckung und Übergangsframes pro Gestaltung
           (Mittelwert über die Sitzungen jeder Generation). Zellen: Anteil
           der 100 Zellen eines $10\times10$-Bildschirmrasters, die mindestens
           ein Label enthalten; Abst.: mittlerer Abstand der Ziele von der
           Bildschirmmitte (normierte Einheiten).}
  \label{tab:design}
  \small
  \setlength{\tabcolsep}{4pt}
  \resizebox{\textwidth}{!}{%
  \begin{tabular}{l c c c c ccc c c}
    \toprule
    Gestaltung & Labels & Länge & Sprünge & Übergang & \multicolumn{3}{c}{Anteil der Labels (\%)} & Zellen & Abst. \\
    \cmidrule(lr){6-8}
               &        & (s)   &         & (\%)     & Mitte & Ränder & Ecken & (\%) & \\
    \midrule
    Standard (Gen 1, 4 Sitz.)          & 493 & 82  & 0  & 0           & 7,1 & 25,0 & 23,0 & 97 & 0,383 \\
    Zufällige Sprünge (Gen 2, 5 Sitz.) & 179 & 30  & 16 & 28,5\textsuperscript{a} & 6,3 & 23,0 & 26,3 & 57 & 0,384 \\
    Eckengewichtet (Gen 4, 9 Sitz.)    & 644 & 107 & 20 & 10,1        & 6,1 & 19,5 & 36,1 & 91 & 0,420 \\
    \bottomrule
  \end{tabular}}\\[1mm]
  \parbox{0.95\linewidth}{\footnotesize \textsuperscript{a}\,Geschätzt:
  Labels innerhalb von \SI{0,5}{\second} nach einem Sprung.}
\end{table}

\paragraph{Beobachtungen.}
\begin{itemize}[nosep]
  \item \textbf{Die Eckengewichtung wirkt wie beabsichtigt:} Gen~4 hat
        \SI{36}{\percent} seiner Labels in den Eckumgebungen, verglichen mit
        \SI{23}{\percent} (Gen~1) und \SI{26}{\percent} (Gen~2), und den
        größten mittleren Abstand von der Mitte.
  \item \textbf{Eine kontinuierliche Bewegung deckt den Bildschirm am
        gleichmäßigsten ab:} Gen~1 erreicht \SI{97}{\percent} der
        Rasterzellen. Die kurzen Gen-2-Aufnahmen (\SI{30}{\second}) decken nur
        \SI{57}{\percent} ab.
  \item \textbf{Übergangsframes:} Mit Sprüngen alle 1,4--2,7 Sekunden und
        ohne Phasenprotokollierung fällt bis zu ein Drittel der Gen-2-Labels in
        einen Übergang; Gen~4 hat \SI{10}{\percent}, die exakt protokolliert
        sind und entfernt werden können (\cref{sec:cleaning}).
\end{itemize}

\subsection{Leistung der Personalisierung (E3)}
\label{sec:design_performance}

Der zufällige Split begünstigt Gestaltungen mit langsamer, kontinuierlicher
Bewegung, weil benachbarte, fast identische Frames in Trainings- und
Testdaten landen. Deshalb wurden alle 18 Sitzungen zusätzlich mit dem
sequenziellen Split ausgeführt (\cref{sec:loading}; Basismodell
\code{norm\_subject}, Seed 42). \Cref{tab:design_perf} vergleicht beide Splits
pro Gestaltung, und \cref{tab:region_errors} zeigt die Fehler pro
Bildschirmregion, berechnet aus den gespeicherten Testvorhersagen.

\begin{table}[htbp]
  \centering
  \caption{Testfehler pro Gestaltung mit zufälligem und sequenziellem Split
           (Mittelwert über die Sitzungen; Basismodell \code{norm\_subject}).
           Konst.: Konstant-Baseline der jeweiligen Testframes; Ep.~0: Anzahl
           der Sitzungen, in denen die beste Epoche~0 war (keine Verbesserung
           gegenüber dem Basismodell); $n_{\text{Test}}$: mittlere Anzahl der
           Testframes.}
  \label{tab:design_perf}
  \footnotesize
  \setlength{\tabcolsep}{4pt}
  \begin{tabular}{l cc cccc}
    \toprule
    & \multicolumn{2}{c}{Zufälliger Split} & \multicolumn{4}{c}{Sequenzieller Split} \\
    \cmidrule(lr){2-3}\cmidrule(lr){4-7}
    Gestaltung & Test & Konst. & Test & Konst. & Ep.~0 & $n_{\text{Test}}$ \\
    \midrule
    Standard (Gen 1)          & 15,29 & 29,22 & 21,52 & 28,70 & 1 / 4 & 69 \\
    Zufällige Sprünge (Gen 2) & 23,45 & 29,44 & 32,80 & 34,39 & 2 / 5 & 22 \\
    Eckengewichtet (Gen 4)    & 19,76 & 31,31 & 28,07 & 31,82 & 5 / 9 & 91 \\
    \midrule
    Alle (18)                 & 19,79 & 30,32 & 27,93 & 31,84\textsuperscript{a} & 8 / 18 & 72 \\
    \bottomrule
  \end{tabular}\\[1mm]
  \parbox{0.9\linewidth}{\footnotesize \textsuperscript{a}\,Mittelwert der
  Sitzungswerte des sequenziellen Splits.}
\end{table}

\begin{table}[htbp]
  \centering
  \caption{Testfehler pro Bildschirmregion (\cref{sec:partition}, $r=0,15$),
           Mittelwert über die Sitzungen einer Gestaltung; in Klammern die
           Gesamtzahl der Testframes in der Region. Basismodell
           \code{norm\_subject}.}
  \label{tab:region_errors}
  \footnotesize
  \setlength{\tabcolsep}{3.5pt}
  \resizebox{\textwidth}{!}{%
  \begin{tabular}{l cccc cccc}
    \toprule
    & \multicolumn{4}{c}{Zufälliger Split} & \multicolumn{4}{c}{Sequenzieller Split} \\
    \cmidrule(lr){2-5}\cmidrule(lr){6-9}
    Gestaltung & Alle & Mitte & Ränder & Ecken & Alle & Mitte & Ränder & Ecken \\
    \midrule
    Standard (Gen 1)          & 15,3 & 17,7 (20) & 13,5 (67) & 16,7 (83)  & 21,5 & 11,9 (25) & 20,9 (64)  & 28,9 (52) \\
    Zufällige Sprünge (Gen 2) & 23,4 & 13,7 (12) & 21,6 (34) & 26,1 (42)  & 32,8 & 6,6 (2)\textsuperscript{a} & 29,8 (15) & 38,6 (45) \\
    Eckengewichtet (Gen 4)    & 19,8 & 11,5 (45) & 14,9 (175) & 26,5 (299) & 28,1 & 9,0 (60) & 23,9 (136) & 37,1 (289) \\
    \midrule
    Alle (18)                 & 19,8 & 13,5 (77) & 16,4 (276) & 24,2 (424) & 27,9 & 9,7 (87) & 24,3 (215) & 36,1 (386) \\
    \bottomrule
  \end{tabular}}\\[1mm]
  \parbox{0.95\linewidth}{\footnotesize \textsuperscript{a}\,Nur eine Sitzung
  mit zwei Frames in der Mitte; nicht interpretierbar.}
\end{table}

\paragraph{Sequenzieller Split.} Ohne Beinahe-Duplikate steigt der Fehler in
jeder Gestaltung, im Mittel von \SI{19,8}{\percent} auf \SI{27,9}{\percent}
(Wilcoxon $p=4\cdot10^{-5}$). Das personalisierte Modell ist immer noch besser
als die Konstant-Baseline (Wilcoxon $p=0,002$), die Verringerung schrumpft
aber von \SI{35}{\percent} auf etwa \SI{12}{\percent}. In 8 von 18 Sitzungen
(02, 05, 08, 23--27) ist die beste Epoche~0: Das Training auf den ersten
\SI{70}{\percent} der Aufnahme verbessert den Fehler auf dem
Validierungsblock nicht, das angegebene Ergebnis ist also das unveränderte
Basismodell. Das geschieht vor allem in Gen~4 mit unruhigem Hintergrund
(23--27) und in den kurzen Gen-2-Aufnahmen, deren Validierungs- und
Testblöcke nur 17--27 Frames enthalten. Der zufällige Split aus
\cref{sec:sessions} überschätzt daher den Nutzen der Personalisierung bei
schwierigen Sitzungen stark.

\paragraph{Gestaltungen.} Bei beiden Splits hat das Standardvideo (Gen~1) den
niedrigsten und die zufälligen Sprünge (Gen~2) den höchsten Fehler; auch mit
dem sequenziellen Split sind die Unterschiede nicht signifikant
(Kruskal-Wallis $p=0,06$). Die Rangfolge lässt sich durch die Abdeckung
erklären (\cref{tab:design}): Gen~1 deckt in jedem Teil der Aufnahme
\SI{97}{\percent} des Bildschirms ab, während bei den Sprung-Gestaltungen ein
zeitlicher Block nur die in dieser Zeit gezeigten Ziele abdeckt. Die kurzen
Gen-2-Aufnahmen leiden am stärksten darunter. Die Gestaltungen sind weiterhin
mit Beleuchtung, Hintergrund und Brille vermengt, die Rangfolge ist also eine
Tendenz und kein Beweis.

\paragraph{Regionen.} Die Ecken sind in jeder Gestaltung und bei beiden
Splits die schwierigste Region (\SI{24}{\percent} zufällig,
\SI{36}{\percent} sequenziell), die Mitte die einfachste
(\SI{10}{}--\SI{14}{\percent}). Die eckengewichtete Gestaltung senkt den
Eckfehler \emph{nicht}: Gen~4 hat \SI{26,5}{\percent} / \SI{37,1}{\percent}
in den Ecken, verglichen mit \SI{16,7}{\percent} / \SI{28,9}{\percent} beim
Standardvideo. Mehr Ecklabels allein helfen nicht, wenn der Blick in den
Ecken schwerer zu schätzen ist (Kopfdrehung, teilweise geschlossene Augen,
Kamerageometrie) und wenn die Ecklabels unter schlechteren Bedingungen
aufgenommen wurden (Gen-4-Sitzungen 23--27). Die äußersten Ecken sind außerdem
die Punkte, an denen die Sigmoid-Ausgabe nahe 0 oder 1 liegen muss.

\paragraph{Folgerung für die Gestaltung.} Für eine kurze Kalibrierung liefert
mit den vorhandenen Daten eine kontinuierliche Zielbewegung, die den ganzen
Bildschirm abdeckt, das beste Ergebnis. Sprünge erzeugen Übergangsframes
(Gen~2: bis zu \SI{28}{\percent} der Labels) und müssen, wenn sie verwendet
werden, protokolliert werden, damit diese Frames entfernt werden können. Die
Eckengewichtung erhöht den Anteil der Ecklabels wie beabsichtigt, senkt aber
nicht den Eckfehler.

% ==============================================================================
\section{Diskussion}
\label{sec:discussion}
% ==============================================================================

\paragraph{Hauptergebnis.}
Auf GazeCapture erreichen alle Konfigurationen der Hyperparameterstudie
praktisch den Fehler eines konstanten Schätzers auf denselben Testdaten: Die
besten Konfigurationen mit \enquote{nur FC} liegen zwischen
\SI{28,32}{\percent} und \SI{28,39}{\percent}, die Konstant-Baseline bei
\SI{28,46}{\percent}. Die räumliche Analyse bestätigt, dass sich die
Vorhersagen in der Bildschirmmitte konzentrieren. Folglich trennt die
Hyperparameterstudie die Konfigurationen vor allem nach ihrer
\emph{Stabilität}: Optionen, die das Training weniger stabil machen (Gauß,
ReLU, keine Aktivierung, Auftauen von Backbone-Schichten, sensitiver Verlust),
sind schlechter, aber keine Option verlässt deutlich die Mittelpunktlösung.
Die Entscheidungen in \cref{sec:training} beruhen daher auf Stabilität,
Konvergenz, Aufwand und methodischen Argumenten, wie in jedem Unterabschnitt
angegeben. Einzelne Notebook-Läufe auf \code{norm\_random}
(\cref{fig:lr_curves,fig:normalized_targets}) erreichten etwa
\SI{23,6}{\percent} gegenüber \SI{26,8}{\percent} für die Konstant-Baseline;
ein kleiner Lerneffekt existiert also, ist aber in den Mittelwerten über fünf
Läufe nicht sichtbar. Die Basismodelle der Personalisierung bestätigen das:
Ihr Fehler auf den persönlichen Sitzungen (\SI{30,6}{\percent}) entspricht
der Konstant-Baseline.

\paragraph{Mögliche Ursachen.} Das neue Protokoll schließt bereits zwei
Erklärungen aus, die für die früheren Läufe plausibel waren: Die
Trainingsteilmenge enthält jetzt 179 statt fünf Personen, und die Testmenge
wird nicht mehr für die Modellauswahl verwendet; der Kollaps bleibt bestehen.
Die wahrscheinlichsten verbleibenden Ursachen sind:
(i) Es werden vollständige Kamerabilder verwendet, in denen die Augen nur
einen kleinen Teil des Bildes einnehmen, sodass die eingefrorenen
ImageNet-Merkmale wenig Blickinformation tragen;
(ii) \num{1000} Frames von 179 Personen sind nur etwa sechs Frames pro Person,
zu wenig, um eine personenunabhängige Abbildung zu lernen; das Auftauen
weiterer Schichten mit so wenigen Daten macht das Training instabil statt
besser (\cref{sec:layers}).
Gesichts- oder Augenausschnitte als Eingabe und größere Trainingsmengen sind
die nächsten zu prüfenden Schritte.

\paragraph{Personalisierung.}
Mit demselben eingefrorenen Backbone senkt ein Kopf, der auf einigen hundert
Frames einer Person trainiert wurde, den Fehler gegenüber der
Konstant-Baseline um etwa ein Drittel (zufälliger Split), und das in etwa
zwei Sekunden auf einer CPU. Die Wahl des Basismodells spielt keine Rolle.
Das stützt die zentrale Idee dieser Arbeit, die dateneffiziente
Personalisierung eines vortrainierten Modells. Zwei Einschränkungen sind
wichtig: (i) Der zufällige Split ist optimistisch; mit einem zeitlichen Split
sinkt die Verringerung im Mittel auf etwa \SI{12}{\percent} und in schwierigen
Sitzungen auf null, während gute Sitzungen den größten Teil des Gewinns
behalten und sich eine Kalibrierung sogar auf eine zweite Aufnahme überträgt
(\SI{47}{\percent} Verringerung, \cref{sec:loading}); (ii) die
Aufnahmebedingungen (Hintergrund, Beleuchtung, Brille) beeinflussen das
Ergebnis stärker als die Gestaltung der Sitzung.

\paragraph{Folgen für den Code.} Die Evaluationsmetrik wurde korrigiert,
Baselines werden auf derselben Teilmenge wie das Modell berechnet, zufällige
Teilmengen (\code{--sample-seed}) und eine personendisjunkte
Validierungsmenge (\code{--val-size}) wurden zu \code{exp.py} hinzugefügt, und
vollständiges Fine-Tuning (\code{fc\_resnet}) kann mit \code{--include-resnet}
einbezogen werden. Die Personalisierungsläufe sind reproduzierbar (feste
Seeds, eingefrorene BatchNorm, Merkmals-Cache, gespeicherte Vorhersagen,
\code{personal\_exp.py}).

% ==============================================================================
\section{Zusammenfassung der experimentellen Ergebnisse}
\label{sec:summary}
% ==============================================================================

\Cref{tab:summary} fasst alle Experimente zusammen. \enquote{Signifikant}
bedeutet $p<0,05$ im jeweils im Abschnitt angegebenen Test. Alle Fehler sind
Diagonalfehler (\cref{eq:ediag}) in \si{\percent}. MAE und MSE wurden in
diesen Experimenten nicht protokolliert; sie werden daher nicht angegeben,
und der Diagonalfehler (ein skalierter euklidischer RMSE) ist das einzige
Genauigkeitsmaß, das für alle Experimente verfügbar und vergleichbar ist.

\begin{table}[htbp]
  \centering
  \caption{Zusammenfassung der experimentellen Ergebnisse und der gewählten
           Konfiguration. GazeCapture: \code{norm\_subject},
           \num{1000}/\num{200} zufällige Frames, fünf Läufe,
           Konstant-Baseline \SI{28,46}{\percent}; Zeit = mittlere Dauer eines
           Trainingslaufs in Minuten (Kopf auf zwischengespeicherten
           Merkmalen). Personalisierung: Mittelwert über die Sitzungen, Zeit =
           ein Sitzungslauf in Sekunden (\cref{sec:session_time}).}
  \label{tab:summary}
  \footnotesize
  \setlength{\tabcolsep}{4pt}
  \begin{tabularx}{\textwidth}{>{\raggedright\arraybackslash}p{2.3cm} >{\raggedright\arraybackslash}p{2.2cm} r r Y}
    \toprule
    Experiment & Gewählt & Fehler (\%) & Zeit & Ergebnis und Begründung \\
    \midrule
    Baselines        & ---                 & 28,46 & --- & Konstant-Baseline auf denselben Testdaten; Zufalls-Baseline 40,36 (\cref{sec:baselines}) \\
    Zieldarstellung  & normiert $[0,1]^2$  & --- & --- & geräteunabhängig; Metrik wohldefiniert (\cref{sec:normalization}) \\
    Batchgröße       & 32                  & 28,39 & 0,019\,min & kein signifikanter Unterschied (\cref{sec:batch}) \\
    Lernrate         & $10^{-4}$           & 28,37 & 0,020\,min & kein signifikanter Unterschied; beste Konvergenz in den Kurven (\cref{sec:lr}) \\
    Trainierbare Schichten & nur FC        & 28,36 & 0,015\,min & niedrigster und stabilster Fehler; FC + Layer 4: 9,2\,min, voll: 15,4\,min (\cref{sec:layers}) \\
    Aktivierung      & Sigmoid             & \textbf{28,36} & 0,022\,min & \textbf{signifikant} besser als Gauß, ReLU, Keine (\cref{sec:activation}) \\
    Verlustfunktion  & L1                  & \textbf{28,32} & 0,011\,min & \textbf{signifikant} besser als der sensitive Verlust (\cref{sec:loss_results}) \\
    Datensatzgröße   & \num{1000}/\num{200} & 28,39 & 0,018\,min & größere Menge nur wegen anderer Testmenge besser (\cref{sec:size}) \\
    Datensplit       & \code{norm\_subject} & 28,39 & 0,020\,min & einziger Split, der die Generalisierung auf unbekannte Personen misst (\cref{sec:split}) \\
    \midrule
    Personalisierung (E1) & nur FC, 18~Sitzungen & \textbf{19,79} & 1,9\,s & \textbf{signifikant} besser als Konstant-Baseline (30,32) und Basismodell (30,57), $-$\SI{35}{\percent} (\cref{sec:sessions_results}) \\
    Basismodell (E0/E1) & \code{norm\_subject} & 19,79 & 1,9\,s & gleich wie \code{norm\_random} (19,59; 2,4\,s; $p=0,20$); gewählt, weil es keine Information über die Testpersonen enthält \\
    Label-Bereinigung (E1b) & WAIT-Frames entfernen & 15,88 & --- & vs.\ 16,72 auf denselben MOVE-Testframes; $-0,84$, nicht signifikant; WAIT-Frames selbst 37,7 (\cref{sec:cleaning}) \\
    Laden der Daten (E2) & sequenziell / zwei Videos & 15,55 / 17,24 & 2--5\,s & zufällig 13,85 ist optimistisch; zwei Videos signifikant schlechter als sequenziell; Kalibrierung überträgt sich auf ein neues Video ($-$\SI{47}{\percent}) (\cref{sec:loading}) \\
    Sitzungsgestaltung (E3) & Standardvideo & 15,29 / 21,52 & --- & zufälliger / sequenzieller Split; zufällige Sprünge 23,45 / 32,80, eckengewichtet 19,76 / 28,07; nicht signifikant und vermengt; Eckengewichtung senkt den Eckfehler nicht (\cref{sec:design_performance}) \\
    \bottomrule
  \end{tabularx}
\end{table}

\renewcommand{\bibname}{Literaturverzeichnis}
\bibliographystyle{unsrtnat}
\bibliography{references}

\end{document}
